\documentclass[12pt,a4paper,fontset=fandol]{ctexart}
\usepackage[top=2.5cm,bottom=2.5cm,left=2.6cm,right=2.6cm]{geometry}
\usepackage{amsmath,amssymb,graphicx,booktabs,longtable,array,float}
\usepackage{enumitem,fancyhdr,titlesec,setspace,xcolor,listings}
\usepackage{tikz}
\usetikzlibrary{arrows.meta,positioning,shapes.geometric}
\usepackage{hyperref}
\hypersetup{colorlinks=true,linkcolor=blue!60!black,urlcolor=blue!60!black}
\onehalfspacing
\setlength{\headheight}{16pt}
\setlength{\emergencystretch}{2em}
\pagestyle{fancy}\fancyhf{}
\fancyhead[L]{\small 海淀区 Demo 实施与验收}
\fancyhead[R]{\small 中国城市建设研究院有限公司}
\fancyfoot[C]{\thepage}
\titleformat{\section}{\Large\bfseries}{\thesection}{1em}{}
\titleformat{\subsection}{\large\bfseries}{\thesubsection}{1em}{}
\setlist{nosep,leftmargin=2em}
\lstset{basicstyle=\small\ttfamily,breaklines=true,frame=single,
 backgroundcolor=\color{gray!5},showstringspaces=false,columns=fullflexible}
\newcommand{\file}[1]{\nolinkurl{#1}}
\newcolumntype{P}[1]{>{\raggedright\arraybackslash}p{#1}}
\title{\textbf{GEP 智能核算详细实施方案}\\[0.5em]\Large AlphaEarth海淀区2024年度全域GEP核算}
\author{伍炜杰}
\date{2026年9月21日修订}
\begin{document}
\maketitle
\thispagestyle{empty}
\begin{abstract}
本方案选用\textbf{2024年AlphaEarth年度卫星嵌入}，开展\textbf{海淀区2024年度全行政区域生态产品总值（GEP）核算}。该年份具有可核实的年度嵌入目录、海淀统计公报及统计年鉴入口，并可获取同期NPP、ET与质量层，便于把遥感、气象水文和社会经济资料统一到完整自然年。依照GB/T 46869--2025回溯核算2024年，覆盖物质产品、调节服务和文化服务三大类，对22个科目逐项登记适用性、实物量、价值量、来源及审核结论，形成总GEP、全域成果、国标核算表和证据包。

AlphaEarth提供生态分类与参数分区所需特征，年NPP、蒸散、生产投入、游客人时及当地价格等仍由独立数据或经审核模型提供。实施按“年度底座、各项计算、去重及外源扣除、总账汇总、人工修正、复算导出”组织。先完成全类型数据收集、清洗、编目和本地数据包验收，再开展应用开发；每个发布版本完成整理、许可检查和复核后，最后同步至ModelScope数据集仓库。公开资料不等于全部输入已验收；文中明确列出可直接获取的数据、需进一步提取的字段和必须补采或估算的缺口。

本文是实施方案，不代表已经导出海淀AlphaEarth资产或算得实际GEP。数据与参数统一使用2024年度，全年366天，不再沿用跨年半年折算，也不把期末碳储量计入年度服务价值。
\end{abstract}
\newpage\tableofcontents\newpage

\section{全区域、全口径的核算目标}
\subsection{空间与时间边界}
\begin{itemize}
\item \textbf{区域：}海淀区全行政区域（110108），覆盖森林、灌丛、草地、湿地、农田、城市等实际存在的生态系统；建成区不整体排除，其公园绿地、道路绿带、景观与游憩服务等按适用科目核算。
\item \textbf{核算期：}$T=[2024\text{-}01\text{-}01,2025\text{-}01\text{-}01)$，共366天，采用完整自然年度，正式成果统一标注“2024年度”。
\item \textbf{统计单位：}实物量以“t/a、m$^3$/a、人次/a”等表达；货币统一为“元/a或万元/a”。面积、库容和期末存量单列，不与期间流量相加。
\item \textbf{空间单元：}以统一行政边界为总账，年度嵌入瓦片与影像像元用于空间建模；林班、地块、水体、道路、景区和受益房产分别使用适合的对象单元。街镇成果在边界通过核验后生成，区级统计数据不强制伪分配到像元。
\item \textbf{交付范围：}三大类全部适用科目进入总GEP核算。数据不足的科目列补采或有依据的估算任务，不以“仅固碳＋水源涵养”作为最终交付。
\end{itemize}

\subsection{总GEP如何形成}
总GEP由三类生态贡献汇总，并通过来源归属与重复项调整形成区域年度总账。按国标第4.8条和第5章，先编制适用清单，再计算和汇总。设$\mathcal{P}$、$\mathcal{R}$、$\mathcal{C}$分别为经确认的物质、调节、文化科目集合：
\begin{equation}
 GEP_T=\sum_{k\in\mathcal{P}}V_{k,T}+\sum_{k\in\mathcal{R}}V_{k,T}
       +\sum_{k\in\mathcal{C}}V_{k,T}-V_{\mathrm{outside},T}-V_{\mathrm{overlap},T}.
\end{equation}
\textbf{公式解读：总GEP＝物质产品价值＋调节服务价值＋文化服务价值－区域外来源价值－尚未消除的重复价值。}三个求和符号表示把相应类别内已审核的分项价值相加，实物量先按各科目方法换算为价值量，再统一采用2024年度、同一价格基期的元/a汇总。
\begin{itemize}
\item \textbf{物质产品：}本地生态系统对农林渔产品、淡水和生态能源等的贡献。人工经营产品按国标扣除劳动、中间消耗、折旧、生产税净额和正常资本回报，得到生态贡献。
\item \textbf{调节服务：}固碳、水源涵养、土壤保持、空气净化等所产生的服务价值，由各自实物量和适用定价方法形成。
\item \textbf{文化服务：}自然景区旅游康养、公园休闲游憩和生态景观带来的年度服务价值，分别通过使用记录、旅行与时间成本、景观增值归因取得。
\item \textbf{$V_{\mathrm{outside},T}$：区域外来源扣除。}若分项初算混入外调水、外购转售农渔产品等不由海淀生态系统提供的贡献，在此剔除。按来源归属，而非按购买者所在地判断；海淀生产后销往外地的产品不因此自动扣掉。
\item \textbf{$V_{\mathrm{overlap},T}$：重复价值调整。}同一贡献已在两个科目计入时只留一次，例如农产品与作物授粉的国标选项、供水法涵养与淡水产品、同一旅游活动的重复时间/费用。这个符号表示本项目的审计调整账，用于落实国标去重要求。
\end{itemize}
\textbf{只扣一次：}优先在分项里排除外源和重复项；总账两个扣减项只登记尚未扣过的余额，若已全部处理则均为0。每笔扣除带来源、理由、金额、入账位置和唯一ID。人工投入在相应产品价值计算中扣，不再混进这两个调整项重复扣。支持服务不另行货币化加总。

每个科目使用“已核算（观测/模型/估算）、不适用、已核实为零、待补数据”四类状态。缺数据不等于零，不适用必须有地理或业务依据。全部适用科目完成后，发布\textbf{海淀区2024年度GEP总值}；若仍有缺项，展示待完善总表及已完成小计并继续补齐，不能将其冒充全口径总值。采用合理论证的替代估算时可形成“全口径GEP估算值”，同时公开估算科目、贡献占比和敏感性范围。

\subsection{必须交付的成果}
\begin{longtable}{P{2.3cm}P{6cm}P{4.7cm}}
\toprule 成果 & 内容 & 验收依据 \\\midrule\endhead
全类型数据包 & 遥感、统计、台账、调查、价格、方法原件及清洗表，附字典和备份目录 & 数据包准入及哈希、逐字段来源核验 \\
全域空间底座 & 边界、2024年度嵌入、生态分类、对象单元、缺测与修复图层 & 全区面积闭合，区内无遗漏且区外不计入 \\
全口径总账 & 22科目适用登记、三类价值及总GEP、来源与重复调整 & 附录E核算表、三类与总账勾稽 \\
交互系统 & 总览、分项地图/对象表、计算链、审核、版本对比 & 10分钟完整GEP演示 \\
报告与证据 & 按附录D编写报告，附输入、参数、年度一致性规则和运行记录 & 原始证据定位、另一环境复算 \\
\bottomrule
\end{longtable}

\section{年度选择与AlphaEarth数据底座}
\subsection{年度选择}
本案例采用\textbf{2024完整自然年}，将年度嵌入、生态物理量和社会经济资料统一到2024年1月1日至12月31日，共366天。标准版本登记为GB/T 46869--2025，数据年登记为2024，价格基期登记为2024；两者分别存储。

2024年度底座包含Source Cooperative的AlphaEarth年度嵌入，NASA MOD17年度NPP及质量信息、MOD16逐期ET、ERA5-Land气象，以及《2025北京海淀统计年鉴》中的2024统计表。遥感产品先登记年度、资产和覆盖，再完成下载验收；年鉴已从官方压缩包提取原件并逐表复核，以下页码采用书内印刷页码[2,5,9,10,15,16]。

{\small
\begin{longtable}{P{2.4cm}P{5.3cm}P{5.3cm}}
\toprule 表号/页码 & 已核对的2024数据 & GEP用途与仍缺内容 \\\midrule\endhead
5-1，第63页 & 农林牧渔业总产值41289.6万元；农业25954.1、林业15058.4、牧业277.2万元；渔业栏空白 & 核对生产规模；林业产值不等于木材价值，需分产品与投入，不能直接加入GEP \\
5-3、5-4，第66页 & 化肥实物/折纯、农药、农膜；粮食面积667.7 hm$^2$、产量4262.4 t；蔬菜面积639.9 hm$^2$、产量24908.4 t & P1的区级数量约束；投入实物可辅助成本调查，但缺分作物价格和完整成本 \\
5-5，第67页 & 生猪出栏1836头、年末存栏818头、猪肉127.2 t；2024水产品栏空白 & 先判断牧产品适用；水产空白不能当零，需渔业部门确认经营/捕捞与收获 \\
5-9、5-10，第70页 & 乡村旅游493118人次、收入20235.7万元；干鲜果总产3159.5 t，其中鲜果3154.1 t、干果5.4 t & 旅游总体与水果分项；旅游收入非旅行成本，鲜果与干鲜果口径不可混用 \\
5-11，第71页 & 各街镇播种面积/产量，含稻谷911.2 t、玉米2935.1 t、冬小麦324.4 t & 可据街镇、作物建数量表；这三项并非粮食全部构成，不强求其相加等于总粮食 \\
5-12、5-13，第72--73页 & 种业产量/收入；设施农业面积283.1 hm$^2$、收入10392.1万元，按温室/大棚等分组 & 苗木、花卉及设施作物可建分项；花卉产量按脚注的百枝/百盆计，不统一当吨，也不与农业总表重复加总 \\
3-2，第50页 & 城镇单位在岗职工平均工资282780元/年，并分行业 & 文化服务时间价值候选；仅代表对应就业群体，不能无说明用于所有游客和居民 \\
12-2、12-3，第179页 & 星级饭店58家，客房收入219869.0万元，另列星级及经营指标 & 景观服务样本框及量级核对；仍缺入住间夜、生态溢价与非星级住宿 \\
16-1、16-2，第215页 & 工业SO$_2$排放122.8 t、工业烟粉尘18.3 t、工业废水165.1万吨；有空气浓度、降尘与噪声均值 & 提供工业源量与环境背景；全源清单及NO$_x$等另行补采；绿带净降噪用同步对照监测 \\
16-6，第217页 & 绿地14069.69 hm$^2$，公园绿地4702.89 hm$^2$ & 类型面积与公园范围校核；公园绿地为子集，不能相加成总绿地 \\
16-15--17，第224--225页 & 水资源总量22800万m$^3$、水资源配置量53407万m$^3$；另表总用水量31458万m$^3$，有供水来源/使用分类 & 三种量分别留存，不能互换；需向水务确认统计范围、外调水、再生水及入出境路由 \\
16-18，第226页 & 污水处理厂5座，能力15.2万m$^3$/日，再生利用量4501.6万m$^3$ & 登记人工处理与循环利用边界；实际处理水量和污染物去除质量另取运营监测 \\
17-1--6，第231--236页 & 居民收入支出、消费与居住结构，其中有房屋来源构成 & 抽样分层及生活背景；不提供生态景观溢价或逐户住房服务价格 \\
\bottomrule
\end{longtable}}

\textbf{提取时先解决三个问题。}第一，使用指南说明空格可表示不详或无数据，不能自动填零；“***”等符号另有含义。第二，表5-4的“干鲜果品”行列3154.1 t，而表5-10明确干鲜果合计3159.5 t、鲜果3154.1 t，须保留原项并请统计部门确认口径。第三，表16-2把PM$_{2.5}$列为31，而区生态环境局发布30.6[22]，可能与取整有关；保存精度和来源，不静默覆盖。水资源配置量、另一表用水量及公报总量的边界也要单独说明。

\textbf{每条统计数据同时保留原表身份、标准字段和核算用途。}当前本地提取表有43条记录，其中40条为已复核转录值，2条为原表空白，1条为果品口径待确认。已复核转录仅表示数值与原表一致；科目适用、成本、来源归属和去重仍在核算时审核。原件保存到\file{01_raw/statistics/}，标准表保存到\file{02_clean/statistics/yearbook_2024_selected.csv}；区域、年度、表号和指标构成唯一键，另存原件路径、印刷页、PDF页、单位和备注。

\begin{longtable}{P{2.5cm}P{4.9cm}P{5.6cm}}
\toprule 数据组 & 归档与清洗结果 & 智能体调用及应用 \\\midrule\endhead
作物、果品、设施 & 原表5-4、5-10--13；建立产品字典与父子关系，区分播种面积、吨、百盆 & 按产品与年度检索数量证据；向P1/P7传入已审核数量，价格和成本缺项生成补录任务 \\
林业、畜牧、水产 & 原表5-1、5-5；另建经营主体、生产方式和采收批次表 & 优先检查科目适用与本地产出归属；水产空白触发人工选择，行业产值用于交叉校核 \\
污染与水务 & 原表16章；污染物、工业/全源、供用水范围分别编码 & R8先查污染物和达标条件；P5/R1查水源路由；范围冲突时挂起受影响输入 \\
绿地、公园及旅游 & 原表16-6、5-9、12章和3-2；总量与子集单列 & 分类面积校核、对象名录与抽样分层；C1--C3另接客流人时、费用和景观溢价 \\
遥感与气象 & 原件、元数据、QC和处理参数；形成年度COG及分区汇总表 & 分类工具读取嵌入；确定性工具读取NPP、ET、降水等物理量，返回输入版本和空间支持尺度 \\
\bottomrule
\end{longtable}

\textbf{智能体通过数据目录取数，计算工具通过已确认的字段执行公式。}现有CLI的\file{build}命令输出\file{observations}与\file{pending}两个集合，每条含稳定行ID及来源定位。应用开发时将它们接入只读检索工具：用户询问“玉米量从哪里来”，工具按区域110108、年份2024、表5-11和指标检索，返回2935.1 t、印刷第71页及本地原件身份；模型据此解释，数值由结构化结果展示。对于“算农产品GEP”，工具继续检查价格、中间消耗、劳动、折旧、税及资本回报，缺一个必需字段即返回补录表单。原表空白进入\file{pending}，不进入公式数值输入。

\subsection{AlphaEarth通过Source Cooperative免费获取}
首选\href{https://source.coop/tge-labs/aef}{Source Cooperative的Taylor Geospatial年度嵌入数据集}[16]，获取2024海淀相交数据。这里提供公开COG下载和空间索引，\textbf{不需要先在Earth Engine排导出任务}；按索引筛选后在本地裁切，下载服务免费不等于本地存储和处理无需资源。此副本的提供者为Taylor Geospatial，底层数据由Google/Google DeepMind生成，两者均记录在来源表中。

索引直链为：

\url{https://data.source.coop/tge-labs/aef/v1/annual/aef_index.parquet}

\begin{enumerate}
\item 保存索引、README与访问日期，检查实际字段；选\file{year=2024}，将海淀边界转为索引的WGS84后做相交筛选。不能仅按文件名猜位置，不能只取第一个瓦片。
\item 用索引记录定位全部相交COG及配套VRT，先验证少量窗口，再下载必要原始文件或有完整追溯的区内窗口。保留URL、源影像ID、源窗口、字节数、CRS、仿射变换及SHA-256；索引相交不等于有效像元覆盖。
\item \textbf{先修正方向，再裁切。}该数据README指出原COG为bottom-up，使用配套VRT纠正读取方向。最终本地VRT必须引用本地文件；产出不依赖远端链接的标准COG，并以已知坐标和边界叠加验证，不手工翻转数组却漏改地理变换。
\item \textbf{先掩膜，再解量化。}原始有符号8位值$q=-128$是NoData；有效值按$e=\operatorname{sign}(q)(q/127.5)^2$转浮点。运算前转浮点避免整数平方溢出，保留64维A00--A63；记录是否已解量化，禁止二次转换。零向量与缺测不能混同。
\item 以原生米制网格计算区界交叠面积，检查空洞、重复和全部64维；分类/模型训练和推理使用同一编码、处理版本。不同获取渠道不直接混拼，先比较对应区域的投影、有效掩膜和特征。
\end{enumerate}
Earth Engine年度集合与Google GCS作为交叉核验或备用获取渠道[9]。冻结的清单记录实际资产和处理参数，不把“网页可访问”作为全区下载及质量验收完成。CC-BY 4.0署名随数据包和派生成果保留。

\subsection{年度嵌入如何参与GEP核算}
本项目设计采用“年度特征＋专家标签＋国标公式”的路径：
\begin{enumerate}
\item 以2024年判读影像、部门生态类型/林班/土地利用资料构建标签。Dynamic World年度汇总[3]可作候选标签，不能当独立真值；老年份WorldCover只作参考，变化区重新判读。
\item 使用全部嵌入维度训练轻量分类器，例如随机森林；基线与不使用嵌入的光谱分类在同一空间留出样本比较。输出生态类型、概率和歧义掩膜，记录特征、样本、模型及随机种子。空间邻近样本分组切分，避免训练验证泄漏。
\item 将覆盖类别映射到国标生态类型，结合经营/管理边界区分森林与城市绿地、水体与湿地等，不能把某一嵌入维度命名为森林或碳含量。
\item 年度NPP、蒸散、降水、径流等由独立产品或经核验模型提供。嵌入辅助分类和参数分区，但不替代物理量、当地价格、游客人时或生产投入。
\item 10 m嵌入和类别面积可提高空间归属细度；500 m NPP或约9 km气象数据仍按其支持尺度计算和验证，不能因叠加嵌入就宣称所有GEP指标达到10 m精度。
\end{enumerate}
年度嵌入总结全年信息，不能当作12个月独立观测。月尺度覆盖、暴雨、降温需求等过程仍使用2024年相应月/日数据。嵌入分类结果与总GEP之间保留明确的面积、参数及公式依赖。

\subsection{全类型输入清单与获取责任}
以下是正式核算数据包必须具有的内容；来源已公开与字段已获取分别验收，关键专用数据按第3章方法卡补齐。
{\small
\begin{longtable}{P{2.1cm}P{5.5cm}P{5.4cm}}
\toprule 数据组 & 首选来源与整理对象 & 服务范围及主要缺口 \\\midrule\endhead
空间底座 & 确认2024区界/街镇界；Source Cooperative年度嵌入[16]；2024 Sentinel-2与地面标签 & 分类、面积及对象登记；行政身份、标签和区内有效覆盖需验收 \\
生态物理量 & MOD17年NPP[2]、MOD16逐期ET[15]、2024气象观测/ERA5-Land[5]、DEM[6]、土壤资料 & 固碳、水土及气候；NPP/ET质量、呼吸、径流和土壤参数需独立处理 \\
农业/林业/渔业 & 年鉴第5章；农业农村、园林绿化、镇村和经营者2024台账 & 品种量价、种植/采收/养殖批次、物料及劳动资本投入；年鉴不足项逐项申请 \\
水务与能源 & 年鉴16-15--18、北京市水资源公报[12]；水务、电网与设施运营方 & 区内外水路由、河湖调蓄、污水处理、发电/能源实际使用；总量与设计能力不能混用 \\
环境与调节参数 & 年鉴16-1--2、区生态环境局、排放源统计及监测，气象/土壤资料 & 分污染物排放与浓度、净化能力、噪声对照、侵蚀及风场；全源和工业源分开 \\
文化与受益对象 & 年鉴乡村旅游/饭店/工资/住房表；公园景区档案、游客调查、租赁交易样本 & 客流人时、旅行费用、归因和景观溢价；不以综合旅游收入替代 \\
价格与方法 & 2024生产者价格、电价、碳价；投入成本、工程预算与寿命；国标、论文和采购资料 & 价格年、单位、方法适用条件与原文页码；每个估计参数必须有可定位依据 \\
\bottomrule
\end{longtable}}

\subsection{先建全类型本地数据包，再开发核算应用}\label{sec:newdata}
项目第一阶段交付\textbf{独立、可理解、可复算的海淀2024数据包}。不以既有目录或脚本决定技术路线；按本方案的指标和质量要求选择输入，过期、不适用、重复及无法说明来源的材料不进入正式包。确认不再使用的工作副本可从项目活动目录清理，清理清单写明原因；保留构成正式结果证据的原件及必要版本。

统一根目录为\file{data/haidian_gep_2024/}，既放遥感，也放年鉴、生产、水务、能源、污染、文化调查、价格和方法资料。结构如下：
\begin{lstlisting}
haidian_gep_2024/
  README.md
  00_catalog/     # sources, fields, units, gaps, licenses, checksums
  01_raw/
    boundary/ remote_sensing/ climate_soil/
    statistics/ production/ water_energy/ environment/
    culture_surveys/ prices/ methods/
  02_clean/       # same subject groups; typed, validated tables/maps
  03_features/    # ecosystem map, object registry, annual drivers
  04_docs/        # yearbook digest, method cards, QA and change log
  05_release/     # frozen upload payload + inventory (no recursion)
\end{lstlisting}
\textbf{数据目录同时保存文件和字段语义。}来源表至少有数据ID、发布者、标题、来源URL、获取日期、统计年度、区域范围、空间/时间尺度、许可及路径和哈希。字段字典至少有原名、标准名、单位、换算、分母、缺失符号含义、科目及原表页码/单元格。缺口表逐字段记录责任渠道、直接值/估计值、可替代方法和是否阻断核算。

统计与调查表采用可分析的长表，主键为“区域/对象—年度/月份—产品/污染物—指标”；原始Excel/PDF及表头、注释保留。CSV使用UTF-8，含代码字段按字符串，数值/日期明确类型，同时输出Parquet便于程序读取。把空白、未统计、不适用与零分别编码，不在清洗时一概填零。金额同时保留原币值、价格年、万元/元换算；面积统一hm$^2$或m$^2$并保留原单位。

空间数据使用GeoPackage/GeoParquet和带CRS/NoData的GeoTIFF/COG，保留原生网格。统一对象ID连接地块、塘口、水源、场站、公园和价格表，不用文件名猜关联。气象先核累计/均值及时间区间，量纲转换独立成表；遥感先质量筛选，再重投影、裁切和按面积交叠汇总。分类数据用类别适用的重采样；连续特征另定规则，不把所有数据统一插值到10 m后冒充高精度观测。

\subsection{数据整理的操作顺序与阶段验收}
\begin{enumerate}
\item \textbf{列齐再收集：}由22科目反推字段清单，固定2024边界、产品/污染物字典与对象名录；并行收集遥感、部门统计、运营台账、调查、价格和理论资料。每次下载先写获取记录。
\item \textbf{原件落盘：}保存官方压缩包、网页/PDF/表格、原始遥感与配套元数据；远程链接不能代替本地副本。限制资料的备份权限单列，目录中可追溯但不得把受限明细误发布。
\item \textbf{清洗与编目：}逐表选2024列，处理合并表头、类型、单位和缺失，合计与原表回勾；遥感完成坐标、掩膜、解量化、覆盖及质量检查。统计初值与修订值不混用。
\item \textbf{总结与归纳：}编写年鉴摘要、各数据组说明及“指标—字段—来源—公式”矩阵；为估计输入保存参数证据和误差范围，人工复核关键数据及全部货币单位。
\item \textbf{冻结数据基线：}缺项按科目列明，所有必需输入已获得或有通过审核的估计方案、且足以运行时，形成data-v1.0。若尚缺关键量，只能发布准备进度包，不能标为核算就绪。
\item \textbf{准备备份包：}按科目和数据组核对文件清单、总字节数、SHA-256、许可与复算关系，做一次从干净目录加载的演练；通过后再启动正式核算应用开发与基线计算。
\end{enumerate}
数据包验收需要：全区对象归属闭合、所有22科目有输入/适用性结论、公式必需字段可定位、统计合计一致、遥感有效覆盖及补缺可说明、所有估计有方法卡。不能以下载文件数量或压缩包大小作为质量完成率。

\subsection{ModelScope数据集的冻结、上传与恢复}
数据集统一存放在\href{https://modelscope.cn/datasets/WeijieWu/haidian_GEP}{WeijieWu/haidian\_GEP}，\textbf{上传是每次发布的最后一步}。先在本地完成清洗、逐值复核、许可审查及冻结，再同步代码与文档，最后上传冻结目录并回下载验证。远端按不可变版本分目录，后续修正发布新版本，保留历史。

\begin{lstlisting}
repo_type: dataset
repo_id: WeijieWu/haidian_GEP
release: reference-2024-v0.1
scope: reference_only
accounting_ready: false
remote_path: releases/reference-2024-v0.1
\end{lstlisting}
\textbf{参考资料包与核算就绪包分别验收。}首个参考版本只包含已核对的统计事实、待确认事项、字段字典、官方来源链接、方法索引和质量报告。它支持数据查证和输入准备；全区AlphaEarth裁切、生产投入、径流、价格等依赖在缺口表中列明。待全部适用科目输入齐备或估计方案通过审核后，再发布\file{data-v1.0}核算基线。每个版本的数据卡写明实际内容、未覆盖内容及可运行环节。

\textbf{公开包采用逐文件许可清单。}年鉴使用指南载有内部使用及转载限制，原始ZIP和扫描PDF留在本地；公开包保存必要的已复核统计事实、表号页码和官方网站链接。调查、经营明细按授权范围保存；AlphaEarth及派生成果保留署名。对不能再分发的输入，记录本地原件哈希和获取步骤，使用者自行获取后放回指定位置。

\textbf{上传前后的核验对应同一份冻结清单。}CLI从独立发布目录建立临时快照，再次验证快照并仅上传清单文件，关闭SDK本地缓存；按允许后缀扫描、拒绝隐藏凭据文件、符号链接与常见令牌内容，生成每个文件的字节数和SHA-256。\file{verify-release}检查新增、删除和内容变化；任一检查失败时停止上传。SDK固定为\file{modelscope-hub==0.4.3}，仓库类型显式指定dataset，关闭同步删除[28]。同一版本由一个发布进程串行处理，同名远端版本目录存在时拒绝覆盖；网络失败产生的部分发布由维护者核查，重新发布使用新版本号。

上传后从ModelScope重新下载\emph{全部发布文件}，逐一比对本地冻结哈希；将仓库、版本路径、提交标识、文件数、核验时间和清单哈希写入发布回执。回执放在本地目录\file{00_catalog/}，不回写已冻结的发布包。具体命令见第\ref{sec:operations}节。

\section{国标全科目清单与计算路径}
\subsection{22个科目逐项建账}
按国标表1、表2和附录E建账。三大类共7项物质产品、12项调节服务、3项文化服务；下表全部进入系统目录，实际计入规则由适用性与去重审核决定。
{\small
\begin{longtable}{P{0.8cm}P{1.7cm}P{5.6cm}P{4.9cm}}
\toprule 编号 & 科目 & 海淀区实施与实物量输入 & 国标方法、价值与期间处理 \\\midrule\endhead
P1 & 农产品 & 按作物、设施/露地与地块收集2024全年实际收获；不以植被面积代替产量 & A.1；人工种植优先残值法，剔除非生态投入；市场法须明确适用依据 \\
P2 & 林产品 & 初级林产品、采集和采伐台账，核实本地生产 & A.1；按产品与投入分账；期内无采收须有零值依据 \\
P3 & 牧产品 & 核实牧区条件及牧草/放养产品，排除外地购入畜产品 & A.1及表1选项；海淀非牧区的适用性审核后决定，不把销售额当生态贡献 \\
P4 & 渔产品 & 本地水体捕捞/养殖收获及投入；逐水体关联 & A.1；以本期收获和可比价格核算，区外水产品不计 \\
P5 & 淡水资源 & 本地生态系统提供且被使用的水量，区分外调水和再生水人工贡献 & A.1；与A.2水量供给法互斥规则联动，不将全区自来水账单直接计入 \\
P6 & 生态能源 & 本地实际使用的生物质及符合范围的水风光等发电/供能台账 & A.1；按实际期内产出及贡献核算，重复销售/消耗不重复记账 \\
P7 & 其他物质产品 & 花卉、苗木、种子等，检查是否已列农林产品 & A.1；本地初级产品与投入剔除，建立互斥产品编码 \\
R1 & 水源涵养 & 全年12个月降水、地表径流、蒸散及类型面积；外调水单列 & A.4/A.5、A.6；实际期内水量或收支，使用对应年度成本与价格基期 \\
R2 & 土壤保持 & 期间降雨侵蚀力、K、LS、月覆盖C及面积 & A.7--A.10；土壤保持量，泥沙淤积及面源污染两项价值 \\
R3 & 防风固沙 & 期间风场、土壤侵蚀/结皮、糙度、覆盖因子 & A.11--A.12；潜在与实际风蚀差异及治沙替代成本；按全年实际风蚀过程核算 \\
R4 & 海岸防护 & 海淀内陆无海岸，边界和生态清单证据支持“不适用” & A.13--A.14；目录保留，不虚构数值或从总清单删除 \\
R5 & 洪水调蓄 & 植被、水库、湖泊、沼泽、水田分账，使用期内暴雨/水位流量记录 & A.15--A.33；事件累计或容量服务分别处理，不能把库容每月重复加总 \\
R6 & 固碳 & 森林、灌丛、草地、湿地、农田等按适用方法；期初期末生物量/年度固碳速率/年NPP与呼吸 & A.34--A.54；计算期间固定CO$_2$流量，价格基期为2024年；不以期末碳储量作GEP \\
R7 & 氧气提供 & 按A.8.1.3关联A.7.1.4固碳速率路径，准备相应本期固碳量 & A.55--A.56；$32/44$换算和制氧价；与NEP独立方法标签关联 \\
R8 & 空气净化 & SO$_2$、NO$_x$、粉尘等分污染物：浓度、排放、净化能力、生态面积 & A.57--A.59；功能区状态决定路径，量乘治理价，年速率须核验年度及类型适用性 \\
R9 & 水净化 & COD、总氮、总磷等：浓度、入出境负荷、区域排放、污水厂处理量 & A.60--A.62；净化能力或质量平衡路径，扣人工处理贡献 \\
R10 & 局部气候调节 & 期内降温/增湿实际服务日、温差或蒸腾/水面蒸发、能效与电价 & A.63--A.67；只计符合条件的期内日，按2024实际服务日计，不能任意固定天数 \\
R11 & 噪声消减 & 按道路类型抽样，绿带近远侧噪声、长度及有效服务天数 & A.68--A.69；降噪量与隔音墙年维护/折旧成本按服务期计价 \\
R12 & 作物授粉 & 期内作物产量、虫媒依赖系数及物候/开花与收获记录 & A.70--A.71；与P1按表1选项避免重复，保留独立展示及计入开关 \\
C1 & 旅游康养 & 自然景区期内游客人次、来源、旅行时间和直接费用 & A.72--A.75；经调整旅行成本法，不等于全区旅游总收入 \\
C2 & 休闲游憩 & 公园绿地、河湖带的期内游憩人时与当地单位时间工资 & A.76--A.77；人时乘工资，与C1同一活动去重 \\
C3 & 景观增值 & 实际客房间夜、受益自住房面积、服务价格与景观溢价系数 & A.78--A.82；酒店与住房服务价值增值，不把房产成交总价或资产存量加入GEP \\
\bottomrule
\end{longtable}}

\subsection{核算对象、输入与估算顺序}
核算对象是\textbf{海淀生态系统在2024年提供的最终产品和服务}，包括实际获取的物质产品、生态过程产生的调节服务，以及人们使用生态空间获得的文化服务。每项沿“对象登记—实物量—生态贡献价值—区域归属—审核入账”计算。

\begin{itemize}
\item \textbf{物质产品：}登记生产主体、地块或水体、产品品类和实际获取量，再把价格与相同产品边界的人工投入关联起来，形成自然贡献残值。
\item \textbf{调节服务：}登记生态类型、面积及服务过程，按水量平衡、碳收支、侵蚀或污染物通量等计算实物量，按国标方法转换价值。
\item \textbf{文化服务：}登记景区、公园及受益对象，使用客流人时、旅行支出和景观溢价等记录，计算年度使用价值与归因。
\end{itemize}
面积用于确定服务发生的位置与规模，生产和监测记录用于确定年度数量，价格及成本用于形成价值。AlphaEarth参与生态分类与参数分区，独立物理数据及社会经济资料分别提供后续计算输入。

每个科目建立方法卡：服务对象、过程机制、空间单元、原始字段、实物量公式、估算条件、价值公式、参数出处、验证与去重。数据选择顺序为：同口径官方2024分项统计；2024经营/监测台账；有抽样框的统计推断；经过本地校验的过程模型；相似地区参数转移。最后一种必须说明气候、经营方式、物种与年份差异。不能取“全国每公顷GEP”乘海淀面积，也不能用专家随意给出的一个比例把产业收入转成生态价值。

Ouyang等（2020）提供了从最终生态服务、实际使用量到货币价值的核算思路[17]；本方案的科目和最终公式服从GB/T 46869--2025。国家统计局面积和单产调查方法[18]、FAO养殖记录方法[19]用于补充\emph{输入量如何取得}。这些调查估计式是实施方法，不冒充国标新增公式。

\subsection{农产品P1：先编品种清单，再算实际收获}
\textbf{对象和来源。}从年鉴农业章、农业农村部门2024种植业报表、镇村与经营主体台账建立“地块ID—经营者—作物—茬次—播种/收获日—面积—产量”清单。逐项核实粮食、蔬菜、瓜果及其他作物；设施与露地分层。果树的水果按农产品计，不因长在林地就再记林产品；粮食作物、观赏苗木和能源秸秆使用不同产品编码。

\textbf{作物清单由统计表逐级展开，首轮区级核算直接采用已确认的分项产量。}整理人员先录入表5-4的粮食、蔬菜、果品总类，再用5-11展开稻谷、玉米、冬小麦等已列品种，用5-10展开鲜果/干果和苹果、梨、桃、葡萄、柿子、杏、红果、枣、李子、樱桃等品类，最后用5-13补设施类型。该名单表示表中存在的类别；每种是否在2024有产量、是否存在其余品种，以2024列及部门说明逐项确认。种业及苗木表5-12进入P7候选清单，分类结论保留依据。

产品字典保存产品ID、父级ID、原始名称、单位、生产方式、归属科目与汇总角色。对同一空间和年度，选择“父项总量”或“完整且互斥的子项”作为一套计算口径。例如鲜果3154.1 t和干果5.4 t可闭合到3159.5 t；只有粮食三项时保留其余粮食待查。设施蔬菜属于蔬菜分组，建父子关系；作物×设施方式交叉表完整后才分层计价。父项与子项同时展示时标记为总计/明细，计算器只取选定层级。

\textbf{是否识别到具体作物，由计算所需的分辨率决定。}海淀区总量已有可靠表项时，直接采用统计值；只具“蔬菜合计”时，可以在获得同口径数量加权价格和成本样本后以蔬菜组核算。若产品间单价、投入或授粉依赖差异影响结果，则向农业农村部门索取分品种表，或组织有抽样框的分层调查。无法可靠拆分的部分保留“其他/待分解”，由操作者选择补资料或使用经审核的分组估计。

\textbf{遥感识别用于补空间位置、面积和调查分层。}实施顺序为：获取2024地块和种植台账；对齐AlphaEarth年度特征；增加Sentinel-2生长季多时相NDVI及红边等特征以表达物候；用有作物名称、日期和空间位置的样本训练分类器；按独立地块或空间块留出验证，输出混淆矩阵、各类精度和低置信区。遥感能支持区分的类别由样本和验证决定，温室遮挡、细碎混种、相近物候作物归入人工复核。空间分配保留统计总量约束；年鉴已提供数量时，识别结果的作用是定位和校核。

\textbf{量如何求。}有可靠分作物实际产量时直接使用，遥感仅核对种植位置和面积。缺分项产量时按国家统计局“面积×单产”的调查逻辑[18]估算：
\begin{equation}
 \widehat Q_c=\sum_{h,s}A_{chs}\widehat y_{chs},\qquad
 \widehat y_{chs}=\frac{\sum_{u\in S_{chs}}q_u/\pi_u}{\sum_{u\in S_{chs}}a_u/\pi_u}.
\end{equation}
$c$为作物，$h$为镇/设施条件等层，$s$为茬次；$A$为与统计口径一致的该茬面积（hm$^2$），$q_u$为样本实际收获（t），$a_u$为样本面积，$\pi_u$为入样概率。单产单位t/hm$^2$。按播种面积定义单产时，将绝收样地零产量纳入；按收获面积估计时单列绝收面积，不能混用两种分母。复种地块逐茬计，但耕地总面积只登记一次。换算1 hm$^2$=15亩，逐产品统一鲜/干重、水分和计量口径。

2024历史案例优先使用当年收获单、称重记录和测产报告；现在的实地测产只能验证方法，不能冒充2024产量。若仅有区级粮食或蔬菜总量，先作为约束；利用可靠地块面积和分层单产权重分配，保持$\sum\widehat Q_{c,h}=Q_c^{\rm official}$。这种分配是空间归属估计，不是独立测量，需保留未分配区级总账。

\textbf{缺数与校核。}亩产超出本地同类作物合理区间时回查单位、茬次和水分；面积与生产报表、地块影像交叉核对。以年鉴粮食4262.4 t和667.7 hm$^2$反推，综合单产约6.384 t/hm$^2$，仅用于量级检查，不能套给每种作物。只有综合产值而无产量时，不能用单一平均价格反推混合农产品吨数。相似地区单产转移须按设施、灌溉和作物对齐，并对单产及面积分别给出不确定性。

\subsection{林产品P2、牧产品P3和其他物质产品P7}
\textbf{林产品。}向园林绿化部门、国有林场和经营主体取2024采伐、销售/自用、林副产品采收记录。木材用实际出材体积m$^3$/a分树种记录，林副产品分品类以t/a记录，不能把森林蓄积量或全年生物量增加直接当木材产量。优先用检尺/过磅台账；只有采伐面积时，按采伐许可对应小班的采前蓄积密度$b_h$、实际采伐比例$f_h$及出材率$r_h$估计：
\begin{equation}
 \widehat Q_{\rm timber}=\sum_h A_h b_h f_h r_h.
\end{equation}
其中$A_h$为实际作业面积，$b_h$为m$^3$/hm$^2$；$f,r$取作业设计和验收档案，不给固定全国系数。林副产品按实际采收面积乘分层采收单产，只有结实潜力或可采储量不计实际收获。保护地或禁采区域如确无采收，凭管理台账记“核实为零”，不能因为有森林就推算木材采收。林地养护修枝去向另记，薪柴与木材销售不能重复。

\textbf{林产品清单从真实采收活动建立。}先向园林绿化部门核查2024经营主体、采伐许可与作业验收，再逐主体登记木材、竹材及依法实际采集的林副产品；海淀是否存在各品类，以管理和交易记录确认。林业产值15058.4万元先拆分为产品生产、育种育苗、营造林及管护等活动，管护工程费用作为成本或活动背景单列。水果归P1，花卉苗木按P7，实际能源利用的薪柴归P6；同一原料的最终归属带唯一产品批次ID。

缺记录时，数据管理员发出“2024采收证明”模板，字段含主体、许可证/作业单、地点、日期、树种/产品、获取量、单位、去向、单价、投入及经办人。林地所有者确认无采收后可提交零值证据；有许可但无验收时暂停产量，或由核算人员选择作业面积公式并上传小班密度、采伐比例和出材率出处，送专家复核。

\textbf{牧产品。}先依国标表1确认“核算区为牧区”的适用条件。海淀不因存在畜牧销售额就计入全部肉蛋奶；圈养、外地采购与本地放牧分开。经确认适用后，草料用实际割取面积×干草单产，奶蛋用2024日/月生产记录，肉类按本地符合条件的出栏记录及相应活重/胴体重口径计，禁止存栏头数直接乘肉价。年鉴畜牧数据首先用于查清适用性及排除边界，非牧区不适用结论保留证据。

\textbf{牧产品按适用结论建立畜种和产品账。}业务人员先提交海淀是否满足国标牧区条件的书面依据；若不适用，科目保存“不适用＋依据＋审核人”，年鉴猪肉等记录保留为背景数据。适用时再按牛、羊等实际畜种和肉、奶、毛、草料等实际产品展开，并登记饲养方式、放牧范围和本地/外购来源。

肉量首选同口径屠宰或生产报表，只有头数时以当年出栏头数乘抽样平均活重和对应屠宰率估算胴体量；奶量按泌乳动物数、实际泌乳天数和抽样日均产奶量估算；毛按实际剪毛头数乘单次净毛重；草料按实际割取面积、干草单产及含水率统一口径。以上参数来自2024主体记录或同生产方式调查，逐项保存样本、分层、误差及出处。存栏用于“期初＋出生/购入－死亡/售出＝期末”的数量校核。外购饲料、种畜、劳动、设施等投入按产品分摊，肉奶等共用成本保存分摊基准，所有分摊比例合计为1。

\textbf{其他产品。}花卉、苗木、种子逐品类、规格统计实际出圃/收获，单位可为株、盆、kg，不强制换成吨。缺出圃量时用苗圃登记、起苗记录与库存变动核对：当年形成可交付产品量=交付量＋期末产品库存－期初产品库存－外购转售净流入量；在圃未形成可交付产品的多年生苗木另列资产。产值、进货和人工投入依同一产品边界登记，避免同一花卉同时列农业和其他产品。

\subsection{渔产品P4：按经营水体与收获批次核算}
\textbf{对象和原始表。}农业农村/渔业管理部门提供水产品报表、养殖许可及经营名录，经营者提供塘口、鱼种、投苗批次、捕捞日期、过磅、转塘、购入和销售记录。区分养殖与天然捕捞；公园观赏水体、禁止捕捞水体不能自动按“水面×平均产鱼率”核算。

\textbf{水产调查先确认经营与捕捞对象，再取得产量。}数据管理员将年鉴水产空白连同表号页码发给农业农村部门，核对2024经营名录、许可和无经营证明。对确认存在的生产，按鱼、虾、蟹等实际品类及养殖/天然捕捞方式建账，登记塘口面积、养殖方式、起止日期、批次和最终收获。公园湖面按管理用途登记；获得允许生产并实际收获的证据后才进入产量账。销售记录附本地生产/外购转售标记，活重、去头去壳重和加工净重先统一换算依据。

若无主体记录但有主管部门完整区级产量，保留区级值并标记“尚未空间分配”。若区级量与批次量均缺，界面呈现“上传称重/报表”“按面积单产估计”“按投苗收获估计”“暂缓”四项；操作者选择估计后填写参数、来源及上下限，系统仅生成待审核提案。

\textbf{直接量与估算。}优先汇总2024最终出塘或天然捕捞的活重当量；活鱼流通企业的外购转售不计本地生产。缺过磅记录时，分塘口—鱼种—批次使用两条互为核验的路径：
\begin{align}
 \widehat Q_{\rm harvest}&=\sum_{p,b}A_{pb}\widehat y_{pb},\\
 \widehat Q_{\rm harvest}&=\sum_{p,b}N_{pb}^{\rm stock}s_{pb}\bar w_{pb}f_{pb}\times10^{-3}.
\end{align}
第一式的$y$为同养殖方式实际\emph{收获}单产；第二式$N$为入塘尾数、$s$为截至收获时成活率、$\bar w$为平均收获体重（kg/尾，式中$10^{-3}$将kg转t）、$f$为该批次在2024实际收获的比例。多批次、转塘和跨年收获逐批登记，不能把投苗量、年末存鱼与出塘量相加。

FAO养殖管理资料区分收获与扣除投苗重量后的生物净增量[19]。本项目因此同时保存收获量、投苗/外购鱼重量、期末存鱼和饲料投入，解释物质平衡，但不把“净生物增长”与国标产品获取量混为一谈；人工饲料、苗种等非生态贡献在价值残值中处理。仅凭通用成活率不能形成正式估计；有可靠区总量时以其约束塘口结果。\textbf{合成手算示例：}10万尾、成活率0.8、收获体重0.75 kg/尾且当年全部收获，得到60 t；或12 hm$^2$、当年收获单产5 t/hm$^2$，亦为60 t。两式是同一数量的交叉检查，不相加为120 t，示例参数不是海淀实测。

\subsection{淡水P5与生态能源P6：只计本地产生且实际利用的量}
\textbf{淡水。}水务部门2024供用水表、取水计量、入出境水量与水源工程台账为首选。逐供水源和受水对象建立路由账，先从实际使用水量识别本地自然水源份额，剔除外调水、重复循环使用以及处理输配的人工贡献。再生水不直接当新增天然淡水。全区总用水不能代替本地生态供水。若采用国标水源涵养水量供给法，执行不另计淡水产品的选项；无来源分解时记待补，不能以海淀面积占比切分北京市水量。

\textbf{生态能源。}依据国标表1、表2和A.1，建立2024本地区域内薪柴/秸秆使用点、水电、风电、光伏设施名录。发改、农业农村、园林绿化部门提供项目和原料台账，电网与运营方提供投运日期、关口表计、自发自用、上网量和厂用电。物理量与价值分开：发电量不是装机功率，太阳辐射潜力不是已使用的能源，更不能把全区用电量当本地生态能源。
\begin{longtable}{P{2cm}P{6.9cm}P{4.3cm}}
\toprule 类型 & 缺直接计量时的估算 & 参数来源与校核 \\\midrule\endhead
秸秆/薪柴 & $M=\sum_c Q_c r_c f_c^{\rm collect}f_c^{\rm energy}$；按干物质口径的草谷比、可收集比例和实际能源利用比例，必要时转热量$E=M H$（$M$为t，低位热值$H$为MJ/t） & 农村能源台账、称重与含水率；还田、饲料和废弃部分不计能源；与原料销售去重 \\
水电 & $E=\sum_t\rho gq_tH_t\eta_t\Delta t/(3.6\times10^6)$，$q$为m$^3$/s，$H$为m，$\Delta t$为s，结果kWh & 运行期流量、水头、机组效率；生态下泄、不发电流量不计；优先表计 \\
风电 & 按实有机组功率曲线$P(v_t)$逐时累加并扣停机/损耗；或装机kW×经本地校验的等效小时h & 场站风速、功率曲线及2024运行状态；不能用全年8784小时当满发小时 \\
光伏 & $E=P_{\rm STC}(H_{\rm POA}/G_{\rm STC})PR$，其中$H$为kWh/m$^2$，$G_{\rm STC}=1$ kW/m$^2$ & 2024组件面辐照、朝向、遮挡、投运日期与性能比；温度/逆变/停机损失只扣一次 \\
\bottomrule
\end{longtable}
上述是能量守恒与工程换算，用于计量缺口估计。光伏参数定义可参考欧盟JRC PVGIS[20]，但多年平均预测不能直接标为2024实测，需核对数据覆盖并使用当年气象。例如100 kW光伏、年度组件面辐照1400 kWh/m$^2$和性能比0.8，估计112000 kWh；这是合成单位校验，不能替代海淀设施与2024气象记录。没有本地生产设施或实际能源使用时，核实后记零；不能凭太阳辐射图新增一个光伏产业。

\subsection{物质产品计价：从产出与投入计算生态贡献}
优先从分产品增加值开始，按国标A.2：
\begin{equation}
 V_p=(O-IC)-L-D-T_n-rK.
\end{equation}
$O$为当年产出价值，$IC$为中间消耗，$L$为劳动报酬（含有依据估计的家庭劳动），$D$为折旧，$T_n$为生产税净额，$rK$为正常资本回报；与国标$E_w,E_r,E_a,E_e,E_b$一一对应。$O$依据产品实际量和2024生产者价格及符合统计制度的库存变化，不能使用含运输零售利润的终端价格。$IC$包含已消耗的种苗、饲料、肥料、能源和外购服务；资本购置全价不在当年全部扣除，贷款利息与$rK$不能重复扣。先有增加值就不再减一次中间投入。

统计年鉴行业总产值只用作核对。缺主体投入账时，按产品、设施/露地、经营规模抽取成本样本，以生产规模加权推算各投入，依据抽样分层而非统一“生态贡献率”。价格用数量加权均价$\bar p=\sum_mq_mp_m/\sum_mq_m$；仅有外年成本资料时作有依据的价格调整，生产技术差异另审。负残值保留核查，不自动归零，也不为了结果好看更换方法。

\textbf{合成示例，非海淀结果：}某作物产出100万元、中间消耗30万元，则增加值70万元；再减劳动20、折旧8、生产税净额2、资本回报5，生态贡献为35万元。不能先把100万元当GEP再在区总账里减一遍投入。国标A.3市场价值法仅在产品和价格已经符合生态贡献边界时使用，并记录方法选择依据[1,17]。

\subsection{空气净化R8：浓度、排放量、净化量分开取得}
\textbf{机制。}植物气孔吸收气态污染物、叶面阻滞颗粒物，产生的是质量通量。监测站浓度是单位空气体积中的质量；缺少时间、输送与沉降过程，浓度不能乘行政区面积直接变成净化吨数。监测浓度下降也可能来自排放治理和气象变化，不能全部归功于植被。

\textbf{第一步，确定国标分支。}对SO$_2$、NO$_x$和烟粉尘分别登记功能区、污染物及对应统计时长的浓度达标判断。2024区生态环境局公开SO$_2$、NO$_2$、PM$_{10}$年均浓度分别为3、28、54 $\mu$g/m$^3$，PM$_{2.5}$为30.6[22]，可作背景核对，不能据此断定所有功能区、所有时段均达标。NO$_2$浓度与NO$_x$排放、PM$_{10}$/PM$_{2.5}$与烟粉尘排放不同，需明确对应关系，PM$_{2.5}$与包含它的PM$_{10}$不能直接相加。

\textbf{第二步，按A.57/A.58计算。}经审核的达标分支以对应污染物年度排放量作国标估计输入；超标分支以生态类型面积×单位面积净化能力计算。国标A.9.1.2总述按\emph{浓度}判定，而其分项文字出现“排放量超过质量标准”及表号引用疑点；实施登记原文与量纲审查，交领域负责人确认，不比较t/a与$\mu$g/m$^3$。必须区分\emph{国标约定估计量}与直接实测生态去除量。

\begin{longtable}{P{2.4cm}P{6.1cm}P{4.7cm}}
\toprule 原始字段 & 获取渠道与2024口径 & 拿不到时如何处理 \\\midrule\endhead
浓度、站点、功能区 & 海淀/北京市生态环境局公报及月报；向监测部门取2024站点小时/日数据、坐标、有效率与功能区图 & 年均值只支撑相应年均判断；小时沉降模型缺小时数据时不能直接运行 \\
SO$_2$/NO$_x$/烟粉尘排放 & 区生态环境统计的2024分污染物排放清单；排污许可执行报告、重点源台账作核验 & 按本地活动量×匹配行业/燃料产排污系数，治理去除效率单列；移动源、扬尘等另表，不以重点企业合计冒充全区 \\
类型面积与净化能力 & 分类底图及地面核查；国标附录C或同气候、树种、叶面积条件的实测论文 & 每个系数登记污染物、年单位、冠层/土地面积分母、季节及出处；缺依据不借用异地统一系数 \\
治理单价 & 按A.9.2.3取得适用定价/征收依据；当地治理项目成本及处理量用于量级复核 & 税额、污染当量和元/t区别登记；不同污染物共同治理费用不各分摊100\% \\
\bottomrule
\end{longtable}
排放缺口估计可按生态环境部产排污系数体系[21]：$\widehat E_i=\sum_l X_l EF_{il}(1-\eta_{il})$，仅当$EF$为产污系数时再乘去除率；若系数已经是排放系数，不重复扣治理。活动量必须是海淀2024实际燃料消耗、产品产量或行驶活动；参数按行业/工艺匹配，不能由遥感凭空生成。

\textbf{第三步，独立物理复核。}参照Nowak等（2006）的沉降机制[23]，以本地浓度、气象、叶面积/叶季和阻力参数计算
\begin{equation}
 \widehat Q_i^{\rm dep}=10^{-12}\sum_{j,t}C_{ijt}v_{d,ijt}A_j\Delta t-Q_i^{\rm resusp},
\end{equation}
这里$C$为$\mu$g/m$^3$，$v_d$为m/s，$A$为模型定义的有效投影面积m$^2$，$\Delta t$为s，结果t；$Q^{\rm resusp}$为同单位再悬浮修正。所用沉降速度若已含冠层/叶面积效应，不再重复乘LAI。此路线是独立校验或经审定的替代估算，不悄悄取代国标分支，更不直接移用美国城市的系数。最终按污染物$V_{R8}=\sum_iQ_i c_i$计价并保留分项。

\subsection{其他调节服务：原始输入、模型与验证闭环}
以下与后文国标公式约束配合使用。公开数据负责基础驱动，部门台账负责实际服务与受益边界；每个模型都要有独立校核，不能仅展示一张彩色地图。
{\small
\begin{longtable}{P{1.8cm}P{5.6cm}P{5.8cm}}
\toprule 科目 & 必须准备的数据与渠道 & 具体计算、缺口估计与校核 \\\midrule\endhead
水源涵养R1 & 2024站点/ERA5-Land降水、MOD16实际ET、地表径流及类型面积；水务提供径流/供水来源 & 水量平衡逐月汇总；无地表径流时可用经本地率定的降水径流系数或水文模型，并显式作为估计。InVEST年产水$P-AET$不是直接的涵养量，不能不减相应径流就改名[24]；按流域边界建模后归属海淀 \\
土壤保持R2 & 气象降雨强度或适用的侵蚀力模型；DEM；土壤质地、有机质；2024覆盖；水土保持监测与泥沙资料 & 潜在侵蚀$RKLS$减实际侵蚀$RKLSC$，逐单元积分；R不能等于年降水量。按国标算保持量，再用输沙/淤积比例算受益；InVEST SDR可复核入河输沙，不能把全部保土量当入库泥沙[25] \\
防风固沙R3 & 2024逐时风速、土壤粒径/含水率、结皮、粗糙度和植被覆盖，气象及土壤部门/产品 & 按A.11参数化潜在和实际风蚀之差；参考点或输沙观测校核。林地存在不等于固沙量大；不得以通用草地参数代替城市路面，零风蚀须有阈值与地表依据 \\
海岸防护R4 & 海淀行政边界及海岸资料 & 内陆不适用，保留判定；不分配北京市或全国海岸价值 \\
洪水调蓄R5 & 水务2024暴雨与水位流量、河湖库容量曲线、闸门运行、湿地/水田面积；DEM辅助 & 按A.15--A.33分植被、库湖、湿地、水田；暴雨过程入流减出流与储量变化闭合。工程设计总库容不等于年度生态调蓄量，容量法按年度价值结构计，不将每场暴雨都乘完整建设费 \\
固碳R6 & 林草小班/样地端点生物量，MOD17年度NPP及QC，碳库/异养呼吸参数；农田投入与秸秆 & 类型分别选国标路径；NPP不是NEP，必须扣异养呼吸或采用适用转换系数。城市/水体缺测不填零，单位和碳/CO$_2$转换显式；用独立样地或可比研究校核量级 \\
氧气R7 & A.7.1.4路径下的固碳速率、面积及制氧价格 & 根据A.55按CO$_2$质量乘32/44，说明与NEP主结果差别。不能把tC直接乘32/44，也不以碳储量计算年释氧 \\
水净化R9 & 水务/生态环境部门：入出境断面流量与COD、TN、TP浓度；工业/生活等排放负荷；污水厂处理量 & 先按功能区和污染物选能力或平衡路径；断面负荷$L=10^{-6}\sum Cq\Delta t$（mg/L、m$^3$/s、s转t），尽量同步流量浓度。平衡量扣人工去除；储量变化、沉积、漏源造成的残差不能自动全算生态净化 \\
气候调节R10 & 2024气温湿度、蒸腾与水面蒸发、土地类型、温湿度需求日、空调能效及电价 & 按A.63--A.67将适用蒸散潜热等转为等效用电；总ET含土壤蒸发，不得全当植被蒸腾。用当地实际气象确定服务时段，核验无重复的降温/增湿需求 \\
噪声消减R11 & 交通道路/绿带长度与宽度、2024或可调整历史近远侧噪声、交通量、隔声墙造价寿命维护 & 按A.68--A.69；相同交通时段、距离与地形条件对照，扣距离衰减和其他屏障差异。dB为对数指标，严格用国标适用形式，不跨路段相加当总响度 \\
作物授粉R12 & P1作物量价表、作物物候、虫媒依赖程度；本地授粉试验或适用品种文献 & 按A.70--A.71分作物用产量、依赖系数及价格计算；风媒作物不套果树系数。已有本地对照试验优先；与P1按表1选项互斥计入，不两次加总 \\
\bottomrule
\end{longtable}}

\subsection{文化服务C1--C3：实际使用者、使用时长和增值归因}
\textbf{旅游康养C1。}文旅、园林与景区提供2024自然景区名录和去重入园人数、客源地、票务及经营档案；全区旅馆营业收入不能替代自然旅游。对不同客源、季节和团队/散客分层，用同口径游客量扩推旅行问卷：$V_{C1}=\sum_jN_j(t_jw_j+c_j)$。逐人保存交通、住宿、门票、旅行时间与工资/机会成本来源；多目的旅行按时间或专门调查归因，往返费用只记一次。只有旅游收入与人数时，缺旅行时间和归因信息，仍不足以完成旅行成本法。

\textbf{休闲游憩C2。}无门票公园以门岗/客流计数、活动台账和分时抽样建立客流总体，按地点、季节、工作日/节假日扩推：$H=\sum_dN_d\bar t_d$，$V_{C2}=He$。$\bar t$来自对应层的平均停留小时，$e$为当地单位时间工资，年度工资转小时必须写明劳动时间分母。手机定位或传感器不能未经校准当人数；家庭收入不能直接代替个人工资。C1同一游客同一活动时段已计旅行时间的部分不得再次计休闲时间。

\textbf{景观增值C3。}从2024酒店实际入住间夜、住房服务价格/租金、受益自住房面积及生态景观可达性建立样本；部门登记和可追溯交易/租赁资料为首选，当前挂牌价只能作另行情景。采用匹配样本或特征价格回归控制区位、交通、学区、楼龄、面积等，将可归因生态景观的价格比例作为溢价系数，按A.78--A.82对\emph{年度住宿/住房服务}计价。不能用房产成交总价计年度流量，也不能把距公园近的全部房价差都解释为生态贡献。缺可信对照与混杂控制时保留估算区间或待补，不任取统一百分比。展开的年度核算式为$V_{C3}=N_{\rm room}p_{\rm room}r_H+A_{\rm home}p_{\rm annual}r_R$：间夜数乘元/间夜，及受益面积乘元/(m$^2$·a)，分别乘无量纲溢价系数；两项均为元/a。

\subsection{论文和采购案例如何转成海淀的工作要求}
\begin{longtable}{P{3.3cm}P{5cm}P{4.9cm}}
\toprule 依据 & 实际可借鉴内容 & 海淀落地要求 \\\midrule\endhead
Ouyang等（2020），PNAS[17] & 三类最终服务、物理量与核算价格、非生态投入区分 & 所有22科目有“量的来源—价值边界—去重”方法卡，外地结果和单价不搬用 \\
国家统计局与FAO[18--19] & 面积/单产抽样，养殖投入、收获和库存记录 & 作物逐茬，水产逐塘逐批；历史台账与抽样权重入证据包 \\
Nowak等（2006）[23] & 污染物沉降依赖浓度、气象和植被条件 & 建立空气净化独立复核链，和国标约定分支分别出表 \\
河南2024采购包9[26] & 招标文件第49页要求明确本地指标、方法、数据来源并形成核算报告 & 对每个公式逐字段列数据需求，交付总账、方法卡与来源矩阵；不能只交一个展示界面 \\
云南2026技术服务采购[27] & 公告要求对调节服务1--2指标比较方法适用性、准确性及本地参数并试算 & 对固碳和水源涵养先做方法对照，解释差异；对照不缩减全GEP交付目标 \\
\bottomrule
\end{longtable}
采购公告证明的是任务及交付要求，不证明某个模型已经验收有效；论文提供方法依据，不免除海淀参数校验。为每个采用的外部方法记录标准/论文页码、公式、适用域、原单位、本地改动和审核结论。

\subsection{水、土、碳等调节服务计算约束}
\textbf{水源涵养：}在水量平衡路径，按月$m$和单元$i$计算
\begin{equation}
 Q_{w,T}=\sum_{m\in T}\sum_i(P_{im}-R_{im}-ET_{im})A_{im}\times10^{-3},
\end{equation}
其中水深为mm、面积为m$^2$。采用地表径流而非总径流、实际蒸散而非潜在蒸散。ERA5-Land月平均累积变量须核对日平均累积语义后乘当月天数及单位转换；蒸发按ET=$-E$换号，不取绝对值。水量供给路径按A.5使用本地用水、净调入水、出入境水量，并执行不再核算淡水物质产品的规则。

\textbf{土壤保持：}按A.7以期间降雨侵蚀力计算$\sum_{m,i}R_{im}K_iL_iS_i(1-C_{im})A_{im}$，面积用hm$^2$且其他因子量纲与之匹配。价值按A.8--A.10分减少淤积和面源污染；容重、淤积比例及污染物百分数单位须显式转换。含水土保持措施P因子的扩展模型与国标基础路径区分版本。

\textbf{固碳：}按国标的数据可得性顺序选择生物量法、固碳速率法、NEP法，并在森林、灌丛、草地、湿地、农田等类型分别登记，不能所有类型套森林系数。生物量法核验年度端点及碳库范围。NEP路径优先用年NPP减年度异养呼吸；采用转换系数时，年NPP $n_i$以gC/m$^2$/a计、类型交叠面积$a_{ik}$以m$^2$计，则
\begin{equation}
 Q_{c,2024}=\frac{44}{12}\sum_{i,k}\alpha_k n_i a_{ik}\times10^{-6}.
\end{equation}
500 m NPP与细分类面积交叠意味着在原像元内采用相同NPP强度，不意味着反演出10 m NPP。年度$\alpha$须核验当地适用性；有年度异养呼吸资料时按A.52处理。NEP法不再乘$(1+\beta)$，后者属于生物量法的适用项。表C.21华北候选分区的森林0.3599、灌丛0.3731、草地0.2565须核实地域和时期，不能默认全区适用。农田按施肥、秸秆和土壤变化处理；期末碳储量可作解释图层，只把年度固碳服务量计价。

\textbf{氧气提供：}作为完整方案中的正式科目实施。按A.8.1.3要求，用A.7.1.4固碳速率路径的期间量按A.55乘32/44，再按A.56计价。若固碳主结果选NEP路径，另存速率法计算及差异解释，不能无提示复用不匹配的方法。年度速率须与类型面积及2024适用条件一致。

\textbf{净化服务：}空气与水净化均按污染物分账，先用功能区对应的\emph{浓度}限值判断路径，不用“排放吨数”和浓度阈值直接比较。水净化质量平衡按A.61计算入境加本地排放，减出境及污水厂人工处理量。原文个别表号或量纲文字不一致时，登记原文、疑点、量纲检查与专家处理结论，不能让LLM自行修正。

\textbf{气候与噪声：}气候调节按A.63--A.67区分降温和增湿需求，只选本期实际发生的服务日并核验需求；噪声消减按A.68--A.69采用符合范围的道路绿带、实测降噪差与有效服务天数。不能因为影像上有绿地就自动生成两项价值。

\section{2024年度核算的时间、单位与价格处理}
\subsection{全年一致性与产品换算}
核算使用2024自然年度，共366天；所有动态输入、年度总量与价格基期均对齐2024。按国标附录的年度单位输出t/a、m$^3$/a、元/a等，不增加半年折算层。静态DEM和土壤允许使用适用年份的资料，但须记录年代及土地变化影响。
\begin{longtable}{P{2.7cm}P{7.7cm}P{2.8cm}}
\toprule 数据 & 年度处理与校验 & 固定产物 \\\midrule\endhead
AlphaEarth & 选2024年度所有相交资产；不把截止日期2025-01-01当作2025产品 & 年度资产清单 \\
NPP & 原始Npp乘$10^{-4}$得到kgC/m$^2$/年[2]；已缩放产品不再乘一次；转gC时再乘1000 & 物理量及QC掩膜 \\
ET & 46个合成期实际累计量按比例因子0.1转换后相加[15]；不得再乘8；核对年末短期和质量位 & 年ET、逐期QC、有效期数 \\
气象水文 & 12个月覆盖完整，日/月均累积语义先核实；2024二月29天；事件驱动模型用实际事件而非年均替代 & 月表、事件与换算记录 \\
统计资料 & 区分资料出版年与数据年，选2024列；修订值与公报初值并存时统一版本并说明 & 年鉴表号、页码、修订记录 \\
容量与存量 & 库容、降噪能力、住房面积用于相应年度服务计价，不把期末资产/碳储量直接加到GEP & 年度流量与存量分表 \\
\bottomrule
\end{longtable}
年ET完整期数不等于质量良好；先按确定的QC规则筛选，再报告有效覆盖及补缺方案。MODIS无效的水体或建成区不能按零NPP/零ET处理，须按适用生态类型和可靠替代路径补齐。年度嵌入的10 m尺度不改变这些产品的原生支持尺度。

\subsection{定价、量纲与去重规则}
价格使用2024实际或有依据的年度代表价格；先有月价和月数量时可逐月计价再汇总，使用年度单价则说明加权规则。不同年份的工程成本统一至明确的2024价格基期并保存调整依据；情景价格独立运行，不混入真实基线。

水源涵养A.6、洪水调蓄A.33及降噪A.69按国标年度运营/维护成本与折旧结构处理，不再乘半年系数。每项在\file{annual_rules.csv}保存原公式、数据单位、换算、价格基期及量纲检查；国标文字或单位有歧义时保存原文及专家处理结论，不由LLM自动改写。

去重采用固定规则登记：农产品与授粉按表1选项选择计入路径；采用水量供给法时不再计淡水物质价值；旅游和游憩的同一活动人时不可重复；住房/酒店景观溢价与旅行费用中重合部分单列归属；水土保持减少的同一污染负荷不能又作为已发生的水体净化量；林产品、其他产品和生物质能源用统一产品ID防止重复登记。不同服务共享影像不等于重复价值，扣除必须指向同一最终生态贡献。


\subsection{质量检查与手算样例}
全域要求是每个空间单元都被归属并完成适用性登记，包含明确的不适用、核实零值及估算；不能通过缩小边界来提高完整度。区内几何缺口和重复覆盖修复后面积守恒误差目标不超过0.5\%。年度嵌入按区内有效面积统计，月度物理驱动另查12个月覆盖；最终适用科目完成率100\%，并分别公布直接观测、模型和替代估算占比。

分类抽查覆盖山地、农田、水体、建成区绿地及边缘，并按面积和生态类型分层；初始60点仅排查明显问题，正式分类精度与模型验证样本量根据类别、空间相关和误差目标设计。文化服务样本另按地点与活动分层，不能沿用遥感样点充当游客调查。

\textbf{合成算术样例，非海淀实测：}一块1 hm$^2$森林本期NPP=100 gC/m$^2$，假设本期适用$\alpha=0.3599$，则本期固碳量1.319633333 tCO$_2$；60元/tCO$_2$情景价值0.0079178万元。碳价改80后为0.0105570667万元，物理量不变。全GEP只增加两者差额，其他科目不变，\textbf{总GEP不会因此整体增加33.333\%}。

\textbf{总账合成样例：}物质产品100、调节服务200、文化服务50万元/a，区域外扣除20、仍未在分项扣除的重复项5，则总GEP=325万元/a。若20万元已在分项净额扣除，系统不得再扣。另测366天边界、全年汇总、缺项禁止冒充零、模型负值保留审查等情形。敏感性按价格、年度模型参数、调查扩推权重分别扰动，幅度标记为实验设置，不冒充统计置信区间。


\section{可执行的开发流水线}
\subsection{任务、输入与落盘产物}
每个阶段都能单独重跑，后续读取文件和元数据，不依赖聊天上下文。
{\small
\begin{longtable}{P{1.1cm}P{3cm}P{5cm}P{3.9cm}}
\toprule 阶段 & 输入 & 实施动作 & 固定输出 \\\midrule\endhead
S0 & 边界、核算期、22科目 & 审核范围、适用性及许可，生成总GEP配置 & \file{demo.yaml}、边界版本 \\
S1 & 产品及访问账号 & 下载或导出；登记请求、资产与哈希 & \file{data_manifest.csv}、原始文件 \\
S2 & 原始数据 & 云掩膜、时间聚合、单位转换、几何交叠、覆盖检查 & 处理后COG/GeoPackage、\file{qa.json} \\
S3 & 国标与价格原文 & 人工登记公式、参数来源和适用条件 & \file{formulas.json}、\file{parameters.csv} \\
S4 & 三类数据及参数版本 & 并行计算适用科目，校核区域外和重复调整后汇总GEP & \file{results.csv}、\file{chain.json}、日志 \\
S5 & 基线结果与修正请求 & 审核差异，生成子运行，重算受影响节点 & 新运行目录、\file{audit.jsonl}、对比表 \\
S6 & 已完成运行 & 生成地图、报告与证据清单；复核后冻结 & \file{report.pdf}、证据ZIP、\file{checksums.sha256} \\
\bottomrule
\end{longtable}}

\section{本地操作手册：获取、整理与同步}\label{sec:operations}
\subsection{代码仓库、环境与目录初始化}
代码与运行环境存放在\href{https://github.com/go-bananas-wwj/haidian-gep}{go-bananas-wwj/haidian-gep}，数据存放在仓库外的独立目录。仓库首版提供数据表校验、目录构建、公开资料下载、缺数审核示例和发布校验CLI；全22科目引擎、地图界面及智能体工具接入按后续PRD实现。安装者先获得私有仓库读取权限，通过本机Git凭据管理器登录；访问令牌放环境变量或凭据管理器。

\begin{lstlisting}[language=bash]
git clone https://github.com/go-bananas-wwj/haidian-gep.git
cd haidian-gep
python3.12 -m venv .venv
source .venv/bin/activate
python -m pip install .
python -m unittest discover -s tests -v
export GEP_DATA_ROOT="$PWD/../data/haidian_gep_2024"
mkdir -p "$GEP_DATA_ROOT"/{00_catalog,01_raw,02_clean,03_features,04_docs,05_release}
python -m gep --help
\end{lstlisting}
基础CLI只依赖Python标准库，发布命令另装\file{requirements-publish.txt}。运行环境使用Python 3.11及以上；本次实际测试环境为Python 3.12。仓库附Dockerfile供容器构建，数据目录通过卷挂载；容器构建及完整地理计算环境按部署阶段另验收。

\subsection{年鉴下载、原件提取与表格复核}
年鉴从海淀区统计局年度数据入口获取，原始压缩包与提取表分别归档。浏览器打开参考资料[10]的2025年鉴页面，点击附件ZIP下载；CLI的\file{fetch yearbook}固定本次核对的官方附件URL，下载完成记录时间、字节数和SHA-256，已存在原件时停止，避免覆盖。
\begin{lstlisting}[language=bash]
python -m gep fetch yearbook --root "$GEP_DATA_ROOT"
python -m zipfile -t "$GEP_DATA_ROOT/01_raw/statistics/haidian_yearbook_2025.zip"
python -m zipfile -l "$GEP_DATA_ROOT/01_raw/statistics/haidian_yearbook_2025.zip"
python -m zipfile -e "$GEP_DATA_ROOT/01_raw/statistics/haidian_yearbook_2025.zip" "$GEP_DATA_ROOT/01_raw/statistics/extracted"
\end{lstlisting}
ZIP原文件名与项目标准文件名的映射写入来源回执。四个PDF分别对应统计图、使用指南/目录、统计公报及第1--9章、第10--19章。扫描页以原图核对；OCR可辅助定位，复核人员对表头、2024列、行名、脚注及合计逐项确认。前半册PDF页码与印刷页码一致，后半册印刷页码为PDF页码加156；页码字段均从1计数。

清洗表字段固定为区域、年度、表号、印刷页、PDF页、原件路径、指标、数值、单位、状态、备注。空白数值保留空单元格，状态为\file{blank_unresolved}；原文有数但口径待确认则为\file{scope_conflict}。只选2024数据列，字符串单位先保留，随后在字段字典登记换算。每次修订生成新数据版本，原始转录留存。
\begin{lstlisting}[language=bash]
python -m gep check "$GEP_DATA_ROOT/02_clean/statistics/yearbook_2024_selected.csv"
python -m gep build "$GEP_DATA_ROOT/02_clean/statistics/yearbook_2024_selected.csv" --out "$GEP_DATA_ROOT/00_catalog/yearbook_catalog.json"
\end{lstlisting}
成功时校验结果\file{valid=true}，构建结果将观测和待确认记录分开。此处只完成结构、唯一键、年份、单位非空和有限数检查；表格语义、值与原图一致性由人工复核报告提供。输出文件存在时命令停止；新一次构建使用版本化文件名。

\subsection{AlphaEarth下载与海淀裁切操作}
AlphaEarth从Source Cooperative按空间索引获取2024相交瓦片，下载与裁切分两步验收。先下载README和全球索引到本地，再用确认的海淀边界筛选；原始索引可复用，其来源、时间和哈希随所选瓦片清单一起冻结。
\begin{lstlisting}[language=bash]
python -m gep fetch aef-readme --root "$GEP_DATA_ROOT"
python -m gep fetch aef-index --root "$GEP_DATA_ROOT"
\end{lstlisting}
地理处理另建环境，安装并验证GeoPandas、PyArrow、Rasterio及GDAL。执行环境依赖导出与GDAL版本查询并保存结果，使用项目确认边界替换工作边界。本次下载索引共有302466条，字段包括year、path、crs、utm\_zone、geom及WGS84包络；几何列为geom，路径列为path。以下按实际字段筛选，读取后再次检查版本：
\begin{lstlisting}[language=Python]
import geopandas as gpd
idx = gpd.read_parquet("aef_index.parquet")
print(idx.columns.tolist(), idx.crs)
boundary = gpd.read_file("haidian_confirmed.geojson").to_crs(idx.crs)
area = boundary.geometry.union_all()
selected = idx.loc[(idx["year"] == 2024) & idx.intersects(area)]
assert len(selected) > 0
selected.to_parquet("haidian_2024_selected_tiles.parquet")
selected.drop(columns=selected.geometry.name).to_csv(
    "selected_tiles.csv", index=False)
\end{lstlisting}
索引path中的\file{s3://us-west-2.opendata.source.coop/}前缀可换为\file{https://data.source.coop/}，其余对象路径保持一致。将对象扩展名\file{.tiff}替换为\file{.vrt}取得配套VRT，保留原始COG与VRT一一对应，下载到\file{01_raw/remote_sensing/alphaearth_2024/}。逐瓦片执行\file{gdalinfo}确认投影、64波段、NoData和方向；先读一个已知地物窗口，叠加边界检查。配套VRT引用改为对应本地COG后再裁切，避免交付文件继续依赖网络。
\begin{lstlisting}[language=bash]
curl --fail --location --retry 3 --continue-at - --output tile.tif "$AEF_COG_URL"
curl --fail --location --retry 3 --output tile.vrt "$AEF_VRT_URL"
\end{lstlisting}
在同目录运行以下脚本，生成只引用本地COG的VRT；保留原VRT作为来源证据：
\begin{lstlisting}[language=Python]
from pathlib import Path
import xml.etree.ElementTree as ET
tree = ET.parse("tile.vrt")
nodes = tree.findall(".//SourceDataset")
assert len(nodes) == 1
nodes[0].text = str(Path("tile.tif").resolve())
nodes[0].set("relativeToVRT", "0")
tree.write("tile_local.vrt", encoding="UTF-8", xml_declaration=True)
\end{lstlisting}
再检查和裁切，GDAL版本需支持有符号Int8及该Warped VRT：
\begin{lstlisting}[language=bash]
gdalinfo tile_local.vrt
gdalwarp -cutline haidian_confirmed.geojson -crop_to_cutline -r near -dstnodata -128 -of COG tile_local.vrt haidian_tile_q.tif
\end{lstlisting}
\file{AEF_COG_URL}、\file{AEF_VRT_URL}由实际清单填写。裁切脚本逐瓦片处理全部相交数据，输出先保留量化整数；后续按第2.2节掩膜和解量化。跨瓦片统一网格时记录目标投影、像元尺寸、重采样和重叠优先规则。覆盖验收计算海淀内有效面积、缺口及重叠，全部64维一致有效后才标记该输入就绪。本节地理命令需在配置完成的GDAL环境运行，基础CLI的\file{fetch}只负责原件下载。

\subsection{其他遥感、气象与业务输入的获取步骤}
物理量与业务量分别从产品平台和数据持有部门取得，再汇入同一目录和字段字典。每次申请写明110108、2024完整年、时间粒度、空间单元、字段、单位、缺失定义、授权用途及联系人；收到文件保留原附件与回复时间，清洗文件另存。

\begin{longtable}{P{2cm}P{6.2cm}P{4.8cm}}
\toprule 输入 & 操作入口与具体选择 & 落盘和准入 \\\midrule\endhead
NPP与ET & Earth Engine登录已启用项目，打开[2,15]的Code Editor；导入海淀边界，日期筛选2024-01-01至2025-01-01。MOD17选Npp及Npp\_QC；MOD16逐期保留ET及ET\_QC。导出任务完成后从Drive下载GeoTIFF；无GEE账号时使用NASA Earthdata/LP DAAC下载原始产品 & 原生投影与支持尺度保留；NPP乘0.0001；ET乘0.1且按实际合成周期汇总；最后一期为6日。分别保存QC、有效天数、缺测与资产清单 \\
ERA5-Land & CDS注册登录[5]，选2024年1--12月、研究区域及相邻水文建模范围、所需降水/蒸发/径流/温度变量；保存页面生成的请求代码，下载NetCDF或GRIB & 保存请求及变量单位；累积量、月日均和瞬时温度分别换算，按当月天数处理；区域站点用于偏差检查 \\
地形与土地覆盖 & SRTM目录[6]选择模型完整流域范围；2024 Sentinel-2/Dynamic World按区界及年份筛选，保留云/质量和逐期时间 & DEM统一高程与投影，按模型处理洼地；年度类别与独立标签验证。边界和训练样本独立验收 \\
农业林业牧渔 & 先提取年鉴第5章，再向农业农村、园林绿化及主体申请分品类产量、采收、经营方式和成本台账 & 原始表、产品主表、主体/地块/批次、投入和价格分表；共用ID关联。敏感明细访问受控 \\
水与能源 & 向水务申请水源/入出境、逐月流量和工程调度，向运营方/电网申请设施名录、2024计量、投运时间 & 水量路由与设施表分别清洗；本地来源和外来来源可勾稽；计量期间与核算期对齐 \\
空气与噪声 & 年鉴16-1--2及生态环境部门；索取2024站点浓度、源排放、源分类和绿带内外同步监测 & 原始日/小时表与年度表分存；工业/交通/生活源各有覆盖说明，污染物名称和质量单位固定 \\
文化与价格 & 公园景区运营、游客抽样、住宿/租赁样本与2024生产者价格；分别申请人数、停留时间、成本及交易属性 & 保存抽样框、权重、匿名对象ID和适用群体；价格来源含税费、运输及价格年说明 \\
\bottomrule
\end{longtable}

公开平台下载失败时记录URL、日期和错误，保留待获取状态；部门尚未回函时记录申请与负责人。所有缺口进入第\ref{sec:prd}节的人机协同流程，已具备输入的其他科目可继续计算。

\subsection{GitHub、Overleaf与ModelScope同步顺序}
同步按“代码和文档先冻结、数据最后上传”的顺序执行，发布清单串联三个平台的版本。代码仓库只保存脚本、环境、模板和方法说明，大体量原始数据通过数据清单引用。
\begin{enumerate}
\item 在GitHub工作目录运行测试，审阅\file{git diff}与\file{git status --short}，按文件添加并提交；拉取远端更新后普通推送到main。更新涉及数据输入契约时，代码与数据版本同时写入兼容表。
\item 在Overleaf Git目录运行\file{git pull --ff-only origin main}，修改\file{implementation/main.tex}及概述副本。XeLaTeX编译两次，查看目录、表格、公式、流程图和编译日志；确认后提交并\file{git push origin main}。网页项目打开详细版主文件，用XeLaTeX重新编译查看。
\item 将允许公开的文件复制到独立\file{05_release/reference-2024-v0.1/}，填写范围、缺项和许可数据卡，记录GitHub与Overleaf提交号。人工复核通过后生成清单并验证。
\item 安装发布依赖，通过隐藏输入读取ModelScope令牌到当前进程环境，上传后立即清理环境变量。令牌只用于认证；日志、CSV和版本清单均不存令牌。
\end{enumerate}
\begin{lstlisting}[language=bash]
python -m pip install -r requirements-publish.txt
python -m gep seal "$GEP_RELEASE_DIR" --id reference-2024-v0.1
python -m gep verify-release "$GEP_RELEASE_DIR"
read -rsp 'ModelScope token: ' MODELSCOPE_TOKEN
export MODELSCOPE_TOKEN
python -m gep upload "$GEP_RELEASE_DIR" --repo WeijieWu/haidian_GEP
unset MODELSCOPE_TOKEN
\end{lstlisting}
\file{GEP_RELEASE_DIR}指向第3步准备好的绝对目录。成功回执包含回下载核验结果，上传完成和回下载通过分别记录。发布失败时保留本地冻结目录与错误状态；先核查远端部分文件，再用新版本目录重发。代码首版仅支持文本型参考包发布，大型遥感分片和受控明细的发布待各自覆盖及许可验收后扩展清单类型。

\section{产品需求：用户流程、人工选择与算例}\label{sec:prd}
\subsection{用户目标与角色}
用户从一个核算项目进入，逐项完成取数、补数、计算、审核和导出，最终得到带证据的年度GEP结果。数据管理员负责导入和来源，核算人员选择适用方法与参数，复核人员确认结论，查看者浏览报告。以下页面和智能体对话为待实现PRD；已有CLI用于验证数据结构、缺数暂停与审核流程。

项目创建表填写区域、2024年度、边界版本、标准版本和数据包版本。首页展示22科目状态及缺口，已完成小计与完整总值使用不同字段；完整总值只在所有适用科目通过审核、来源与去重调整闭合后生成。核算用户可先运行具备条件的科目，缺数任务保留在同一项目中。

\subsection{从建项目到导出结果的流程图}
\begin{figure}[H]\centering
\begin{tikzpicture}[node distance=0.48cm,
 box/.style={draw,rounded corners,align=center,text width=10.3cm,minimum height=0.75cm,font=\small},
 arr/.style={-{Latex},thick}]
\node[box] (a) {1. 创建海淀2024项目：选边界、标准、数据版本};
\node[box,below=of a] (b) {2. 导入数据包：校验来源、年度、单位与文件哈希};
\node[box,below=of b] (c) {3. 选择科目：展示对象清单、已有输入与缺口};
\node[box,below=of c] (d) {4. 缺数处理：人工填报／选择估算／暂缓该项};
\node[box,below=of d] (e) {5. 预览计算：实物量、价格、投入、公式及来源};
\node[box,below=of e] (f) {6. 人工复核：通过入账；退回时补数并生成新运行};
\node[box,below=of f] (g) {7. 汇总与去重：三类小计、调整账、完整性检查};
\node[box,below=of g] (h) {8. 导出证据报告：冻结版本，数据验收后最后同步备份};
\foreach \u/\v in {a/b,b/c,c/d,d/e,e/f,f/g,g/h}\draw[arr] (\u)--(\v);
\end{tikzpicture}
\caption{核算用户的主流程；暂缓科目保留缺口，其余科目继续}
\end{figure}

\subsection{页面输入、用户动作与系统反馈}
\begin{longtable}{P{2cm}P{4.2cm}P{3.9cm}P{2.9cm}}
\toprule 页面 & 输入与显示 & 用户动作 & 系统输出 \\\midrule\endhead
项目与数据 & 110108、2024、标准版本、边界、发布包及哈希；来源许可摘要 & 选择本地目录/远端冻结版本，导入 & 接入报告与缺项清单；错误定位到文件/字段 \\
科目工作台 & 22科目，适用性、输入齐备率、审核状态；展开到产品/水体/对象 & 选择P1或P4，查看原表；填写适用结论 & 方法卡、所需字段与下一操作；未确认数据保持待补 \\
补数表单 & 缺字段、单位、年度、来源、可选估算及适用条件 & 上传文件/手工填数；选择估算或暂缓 & 参数检查、证据列表、候选结果和不确定性；进入待审核 \\
计算与复核 & 公式、每项数值、来源页、人工投入、去重关系与前后差异 & 核算员提交；复核员通过或退回并填写理由 & 新run、父run与操作审计；通过后仅相关节点重算 \\
总账与导出 & 三类小计、外源扣除、重复调整、未完成项、估计占比 & 查看链路，导出CSV/报告/证据包 & 可复核结果；全项完成时生成总GEP，版本冻结 \\
\bottomrule
\end{longtable}
\textbf{缺数提示给出可执行选择和影响范围。}示例：“P4缺少2024实际收获量，原表5-5为空。补录收获台账；或选择面积×收获单产估计；或暂缓P4。暂缓后继续其他科目，总GEP等待P4完成。”每个估计表单说明参数含义、单位、来源要求、适用条件和敏感性，不预填全国统一系数。界面显示操作者选择过的值，智能体只起草待审提案。

状态依次为\file{waiting_input}、\file{pending_review}、\file{approved}；数据补齐但仍缺价格/成本时转\file{waiting_valuation_inputs}。生产版增加\file{not_applicable}与\file{verified_zero}的证据审核；任何人工修正生成新run，旧运行保持可查。暂停作用于缺数项及其依赖，总账保留完成小计和阻塞项列表。

\subsection{算例一：按统计表计算玉米生态贡献}
用户以年鉴提供的玉米产量启动P1，系统将数量取证与价格成本补录连成一次计算。选择“农产品—玉米”后，数量栏显示2935.1 t，证据为2025年鉴表5-11、印刷第71页；系统提示补入2024生产者均价及同边界成本。

下表仅演示产品交互和算式。\textbf{产量引用年鉴，价格与成本全部为合成教学值，结果不属于海淀实际GEP。}
\begin{longtable}{P{4.8cm}P{3cm}P{5.2cm}}
\toprule 字段 & 示例输入 & 输入及检查方式 \\\midrule\endhead
玉米获取量 & 2935.1 t & 统计表自动带入，点击查看原表证据 \\
生产者均价 & 2000元/t & 用户填报，附价格期间、数量权重、单据；本例标synthetic \\
中间消耗 & 2000000元 & 肥料、种子、能源、外包等同边界合计 \\
劳动报酬 & 1200000元 & 雇佣及家庭劳动估计，记录工日和单价 \\
折旧、生产税净额、资本回报 & 400000、100000、300000元 & 资产、寿命、税补贴及资本基数/回报率分别留证 \\
\bottomrule
\end{longtable}
计算器先得到产出$2935.1\times2000=5870200$元，再得增加值$5870200-2000000=3870200$元，最终残值$3870200-1200000-400000-100000-300000=1870200$元。页面依次展示数量、产出、增加值和生态贡献；复核员确认方法与证据后才可入实际账。本例整体锁定教学模式，用于测试单位与减项顺序。

用户问“这个数怎么来的”，智能体调用结果链，返回年鉴表号、已填成本表、公式版本和本次run；问“价格改成2100会怎样”，系统创建情景run，显示差额293510元及敏感性，其余输入和基线run保持可追溯。情景结果经人工确认后才能作为新的正式参数候选。

\subsection{算例二：水产缺数时由用户选择估算}
用户选择P4后，系统以年鉴空白生成缺数任务，操作者决定取得直接记录、估算或暂缓。选择“面积×收获单产”后，表单要求填写经营塘口、实际养殖面积、鱼种、当年收获单产、来源与参数范围。教学示例输入12 hm$^2$和5 t/hm$^2$，计算60 t；选择“投苗×成活率×体重×当年收获比例”时，10万尾、0.8、0.75 kg/尾和1也得到60 t。

两种方法分别形成提案，用户选一条作为候选，另一条用于交叉核验。系统记录\file{pending_review}，复核员检查真实经营、参数出处、跨年/转塘和外购转售。通过数量审核后仍需补单价及投入，系统显示\file{waiting_valuation_inputs}；60 t只作为已审核数量，总GEP保持空值。若操作者选择暂缓，P4继续\file{waiting_input}，其他具备条件的科目继续。

以下命令可直接在当前CLI运行，生成3个互不覆盖的教学运行文件：
\begin{lstlisting}[language=bash]
python -m gep example --out runs/example-01.json
python -m gep resolve runs/example-01.json --input examples/fish-area-yield.json --out runs/example-02.json
python -m gep approve runs/example-02.json --reviewer reviewer --out runs/example-03.json
\end{lstlisting}
第一份含玉米教学计价与水产缺口，第二份增加60 t待审提案，第三份记录数量审核通过及尚缺估价输入。操作者和复核人使用不同标识；当前CLI记录名字用于演示，生产系统通过登录、角色和服务器审计验证身份。再次运行时更换输出文件名，保留所有父子run。

\subsection{产品验收与异常处理}
交互验收以用户能完成具体任务为准，逐一检查输入、反馈和追溯。导入错误年度文件时定位到行；水产空白触发缺数选择；成活率大于1时阻止提交；未经审核的估计保留提案状态；数量通过而价格缺失时继续等待估价；重复提交同一输出文件时拒绝覆盖。后台正式引擎另外验证父子品类互斥、人工成本完整、全22项适用性及总账勾稽。

网络下载失败可重试并保留来源任务；原表口径冲突进入人工确认；同名发布目录已有远端内容时停止覆盖。审核退回附原因和责任人，修改后创建新版本；发布包只能来自冻结结果，上传与回下载核验完成后记录备份成功。


\section{系统架构与智能体的职责}
\subsection{确定性计算与自然语言交互分工}
\begin{figure}[H]\centering
\begin{tikzpicture}[
 box/.style={draw,rounded corners,align=center,text width=11.8cm,minimum height=1.1cm,font=\small},
 arrow/.style={-{Stealth},thick},node distance=0.6cm]
\node[box,fill=blue!6] (ui) {Streamlit：地图、指标卡、证据查看、人工审核与导出};
\node[box,below=of ui,fill=blue!6] (agent) {Dify：解析意图、定位国标、提出参数修正、解释已计算结果};
\node[box,below=of agent,fill=green!6] (api) {核算API与任务执行器：校验请求、持久化状态、运行固定公式};
\node[box,below=of api,fill=green!6] (data) {本地冻结数据、参数登记、DAG计算链、不可变运行与审核记录};
\draw[arrow] (ui)--(agent);\draw[arrow] (agent)--(api);\draw[arrow] (api)--(data);
\end{tikzpicture}
\caption{全GEP组件及职责；结构化表单可直接调用核算API}
\end{figure}

语言模型负责把“计算海淀区2024年度总GEP，并说明三类贡献”转为白名单请求；不执行任意Python，不自行补写观测数据或价格。核算引擎返回结构化结果和引用，回答中的每个数值都必须对应结果字段。总GEP任务默认加载三类全部适用科目；单指标查询仅用于钻取或专项重算，不改变全口径交付范围。

Dify代码节点只做轻量字段整理；地理空间处理通过HTTP工具调用独立核算服务，避免假设其沙箱内已装GDAL或可访问宿主机数据。首版使用固定工作流，不把ReAct自由规划和平台特定“暂停节点”作为硬依赖。人工审核状态保存在业务服务中，跨页面刷新与LLM超时仍可恢复。

\subsection{API和状态机契约}
\begin{longtable}{P{5.7cm}P{7.7cm}}
\toprule 接口 & 行为及约束 \\\midrule\endhead
\file{POST /runs} & 接收配置版本与幂等键，返回202及run ID；数据不满足准入条件返回缺失清单 \\
\file{GET /runs/{id}} & 返回queued/running/succeeded/failed/blocked、进度与错误原因 \\
\file{GET /runs/{id}/chain} & 返回节点、依赖、产物及可用证据链接 \\
\file{POST /runs/{id}/revisions} & 保存修正草案、理由、操作者和差异；不立即覆盖基线 \\
\file{POST /revisions/{id}/approve} & 人工确认后创建parent run明确的子运行；版本冲突要求重新查看差异 \\
\file{GET /runs/{id}/export} & 仅导出完成运行；未审核包带草稿标记，审核通过后输出冻结包 \\
\bottomrule
\end{longtable}
运行状态与审核状态分开：运行成功不等于科学审核通过。审核状态采用draft、in-review、approved、rejected。修改单位或数据引用也必须重做校验；相同幂等键不得创建重复任务。

\section{“有迹可循”的证据设计}
\subsection{从结果反查原始数据}
点击总GEP，应沿“总值→三大类→科目→实物量与定价→嵌入/监测/统计/调查原件”逐级查看；区域外扣除与去重调整也作为独立节点。固碳节点可以继续追到年NPP、年度异养呼吸或转换系数与边界类型面积。每个节点至少保存：
\begin{enumerate}
\item 节点ID、上游节点ID、运行状态、起止UTC时间；
\item 输入文件路径、SHA-256、资产ID、获取时间、覆盖起止日期、原生分辨率、单位和质量掩膜；
\item 公式ID、代码Git提交、依赖锁文件哈希、容器镜像摘要；
\item 参数原值、单位、来源文档及页码/表行列、估算机制与适用条件、审核人；
\item 输出数值或中间文件引用、单位、有效面积、异常/缺测数量；
\item 人工修正的原值、新值、理由、操作者、审核者及关联新运行。
\end{enumerate}
原始大数组保存为文件并记录哈希，不能仅在JSON中写“已生成10 m分类数据”。遥感质量、数据覆盖率和敏感性分析分别显示，不把来源可靠性随意转换成0.95之类的概率。

\subsection{DAG依赖、不可变版本和复算}
计算链采用有向无环图（DAG），从显式\file{depends_on}生成连线，不能按日志先后把所有节点串成一条线。数据与参数内容哈希组成缓存键，并包括公式、代码和环境版本。若只改碳价，只失效固碳价值、调节类汇总、总GEP及报告节点；NPP与固碳实物量复用，记录复用来源。改年度$\alpha$需重算适用NEP及其下游；改边界必须重做裁切、面积和全部下游。

每次修正生成新run ID和\file{parent_run_id}。旧运行及其证据不被覆盖；重跑报告中比较输入清单、中间数值与容差。相同环境优先比较数值产物哈希；PDF时间戳等可变元数据不作为科学复算一致性的判据。哈希用于发现变化，不宣称单靠哈希即可防止恶意篡改。

\subsection{交付包目录}
\begin{lstlisting}
haidian-demo-v0.1/
  README.md                  # scope, limitations, reproduction
  config.yaml
  data_manifest.csv
  parameters.csv
  indicator_catalog.csv       # all 22 items and applicability
  annual_rules.csv            # annual units, timing and QC
  adjustments.csv             # outside/overlap and booking level
  gep_table_E.csv             # standard Appendix E
  formula_registry.json
  qa.json
  code_version.json          # commit, lock hash, image digest
  runs/baseline/             # result, chain, logs, intermediates
  runs/revision/             # parent reference and changed results
  audit.jsonl
  comparison.csv
  report.pdf
  demo_script.pdf
  screenshots/
  checksums.sha256
\end{lstlisting}
原始数据在许可允许时随包提供；受限数据保存访问路径、校验值和获取要求。只有索引而接收方拿不到原文件时，必须标记“复算依赖未交付”，不能验收为完整离线可复算包。标准原文按授权提供，至少提供条款索引和可定位证据。

\section{总GEP展示界面与10分钟现场脚本}\label{sec:demo}
\subsection{四个页面的明确内容}
\textbf{页面一：全域GEP总览。}顶部显示2024年1月1日至12月31日、366天、边界版本、run ID及审核状态；总卡显示本期总GEP，旁列物质、调节、文化三类小计及区域外/去重调整。展示22科目完成情况和估算占比；缺项时总卡明确“待完善”，不将部分指标包装为全口径总值。

地图提供生态系统底图、各项服务分布及可合理空间归属的GEP分布。水体/道路/景区采用对应对象图层；无法合理分配的区级统计显示为“区级总账”，不制作虚假的逐像元精细价值。西山林地、农田、河湖和城市公园分别演示不同类别，建成区服务不整片空缺。可空间归属值加区级未分配值，应与全区总账一致。数据准备页同时展示第\ref{sec:newdata}节的年度资产覆盖、分类验证、来源版本与哈希和核算阻断项，避免将原始特征展示当作已核算结果。

\textbf{页面二：三类分项与计算链。}左侧选择任一科目，展示实物量、价值、方法、是否计入总GEP及理由；右侧按DAG反查原始嵌入、生态映射、期间参数、统计凭证、调查样本和国标页码。通过“旅游康养”说明非遥感数据的作用，通过“固碳”展示年度嵌入分类与独立生物物理输入的关联。

\textbf{页面三：人工修正与总值重算。}输入“把本次情景碳价由60改80”，先展示差异、影响科目与总账，再由人工确认创建新运行。显示固碳价值的变化、调节类变化和总GEP绝对增量；其他科目复用。另一预设案例修改游客停留时长，展示文化类及总值更新。

\textbf{页面四：国标总表与证据。}按附录E展示全部科目、实物量及单位、价值及单位、三类和GEP合计；追加适用状态、估算标记、年度适用性及扣除行，保持与国标字段映射。报告按附录D组织前言、区域概况、目标原则、方法数据、实物量、价值量、结论及附件，核算基准年位置写清2024自然年度。

\subsection{现场逐步操作与预期结果}
{\small
\begin{longtable}{P{1.3cm}P{5.1cm}P{6.6cm}}
\toprule 时间 & 演示动作 & 必须看见的结果与证据 \\\midrule\endhead
0--1分 & 打开全区本期基线 & 起止日期、总GEP、三类小计、调整项、22科目完成状态 \\
1--3分 & 查看2024年度嵌入覆盖及全域地图 & 边界内完整归属、区外剔除、缺口修复及实际空间尺度；区级总账单列 \\
3--5分 & 分别钻取固碳和旅游康养 & 年NPP或生物量链、国标条款、游客原始凭证；说明遥感与统计如何共同核算 \\
5--7分 & 情景碳价60改80，人工确认 & 新run；固碳价值乘4/3，总GEP仅增加该项差额；物理量与其他科目不变 \\
7--8分 & 查看年度一致性规则及授粉/淡水去重 & 366天、年度输入明细、选项互斥、区域外扣除及调整只记一次 \\
8--9分 & 测试运行撤下一项必需文化服务数据 & 完整性门槛失败、总值待完善；原冻结完整基线不受影响 \\
9--10分 & 导出国标核算表、报告及证据包 & 三类与总账一致，全部适用科目可追溯，包内含复算说明 \\
\bottomrule
\end{longtable}}
60与80是演示用价格情景，不能标注为实际碳市场均价。真实基线必须包含全部适用科目；异常处理演示、历史回放或合成总账样例不能替代最终全域GEP验收。

\subsection{现场稳定性}
演示前冻结数据、模型、镜像、价格与调查权重，完整排练两次。下载与全量计算预先完成；使用本地地图与证据，现场执行可验证的增量重算。LLM不可用时用结构化表单调用同一业务API，历史结果明确标记“历史运行”。

\section{实施排期与全口径验收}
\subsection{八周工作安排及数据依赖}
实施按数据阶段与核算/展示阶段顺序组织，下表按八周作为规划基准。数据包准入是硬性前置条件；专用数据未就绪时顺延后续开发，不在界面中临时拼接来源不明的数据。人员建议包括地理数据、生态/水文模型、社会经济统计调查、后端和前端/智能体各1人，领域与经济核算负责人每周审核；部分角色可兼任，但应相应调整周期。数据申请和历史调查资料获取是关键路径，八周以所需权限和基础资料能按期获得为条件，不承诺凭遥感数据完成缺失统计。

{\small
\begin{longtable}{P{1.2cm}P{2.5cm}P{5.4cm}P{3.7cm}}
\toprule 周次 & 牵头 & 任务 & 验收里程碑 \\\midrule\endhead
第1周 & 领域/数据/统计 & 确认范围、22科目字段和来源；下载年鉴与AlphaEarth，部门台账并行收集 & 本地原件、字段清单与获取时限 \\
第2周 & 地理/模型 & 全类型清洗、统计表提取、单位与边界核对；分类标签和模型输入整理 & 数据字典、来源矩阵及缺口报告 \\
第3周 & 数据/领域/统计 & 完成方法卡、估计输入及质量审查，冻结本地数据包与备份目录 & 数据包就绪；仓库地址由用户后填 \\
第4--5周 & 模型/后端/领域 & 数据准入后实现公式与校验，全科目基线运行，外源与重复审核，三类汇总及总GEP & 总账与附录E表勾稽一致 \\
第6周 & 前端/智能体 & 全GEP页面、分项DAG、修正与局部重算、附录D报告导出 & 完整演示主线可走通 \\
第7周 & 调查/领域 & 独立复核、样本验证、模型参数/价格敏感性，另一环境复算 & 估算范围与模型误差可说明 \\
第8周 & 全体 & 补齐缺项、演练与证据冻结，交付全域总GEP成果 & 所有适用科目及全区验收通过 \\
\bottomrule
\end{longtable}}
第1周明确数据交付日期，第3周审核数据包及缺项；必要时采用有出处、可复核的估算并标记，或延后受影响里程碑。后续缩小为两个指标只能视为内部调试，不能作为本次要求的最终交付。因为核算期已结束，事后新做的调查需说明回忆偏差及时间代表性，不冒充当期现场观测。抽样数据按预设分层和权重计算标准误或重抽样区间；参数转移使用情景范围，两者分别标注，不将随意上下浮动10\%称为统计置信区间。

\subsection{可测量的验收表}
{\small
\begin{longtable}{P{2.3cm}P{7cm}P{3.7cm}}
\toprule 项目 & 判定标准 & 留存证据 \\\midrule\endhead
数据先行 & 开发前完成全类型本地整理和准入；非遥感与遥感同等编目，ModelScope冻结版本可恢复且许可范围明确 & 清单、字典、质量报告、备份目录 \\
范围与期间 & 全海淀边界、2024-01-01至2025-01-01左闭右开、366天；区外剔除及边缘补齐 & 面积表、年度资产清单 \\
科目完整 & 22科目全部有适用结论；所有适用科目完成核算，三类都被评估；不存在以缺测伪装零值 & 适用表、完成率、估算占比 \\
嵌入作用 & 2024资产、64维特征、分类模型、参数及指标可追踪；标签与网格身份、缺测修复和留出验证均有记录 & 模型卡、留出验证、来源图层 \\
年度一致 & 数据和价格对齐2024，年度单位和质量规则可复核；容量、存量与流量分离 & 期间规则及量纲审查表 \\
总账正确 & 三类分项扣调整等于总GEP；外源和重复只扣一次；地理分配值加未分配总账与区合计一致 & 附录E表、调整账、勾稽报告 \\
数值验证 & 固碳手算和325万元总账样例通过；覆盖单位、百分数、缺测、负值、闰年边界 & 检查报告与人工复核 \\
局部重算 & 改碳价只改相应价值及总账，其他科目不变，旧版保留；总值增量与分项一致 & 新旧run、哈希及差异表 \\
证据与复算 & 所有发布值有来源、公式、期间规则和审核；同环境复算相对误差不超过$10^{-6}$ & 证据包、另一环境日志 \\
性能与演示 & 缓存页面目标3秒内、价格重算目标10秒内，各测10次；完整主线连续两次通过 & 中位/最大时延、录像 \\
\bottomrule
\end{longtable}}
性能是待测目标；全量模型、统计核算冷运行单独实测，配置8核CPU、32 GB内存作为初始基线，存储按实际嵌入体积预留。使用已发布年度嵌入，不安排基础模型预训练；分类器训练资源以实际模型测试为准。申报书50并发另做平台验收，不以单用户Demo代替。

\section{部署、版本冻结与文档同步}
\subsection{部署方案修正}
Dify使用官方版本对应的完整Docker Compose配置，记录发布标签与镜像摘要，不使用不存在的单一\file{langgenius/dify:latest}服务替代全套依赖。按官方文档初始化地址为\url{http://localhost/install}，自定义端口以实际配置为准[8]。

核算API和Streamlit单独编排，分别提供服务端口8000与8501。Dify通过配置的内部服务地址调用核算API；容器内\file{localhost}不代表宿主机或另一容器。交付前验证网络连通、请求超时、任务状态持久化及重启恢复。LLM密钥只放环境配置，不写入Git、计算链或导出包。

Python环境采用锁定依赖的镜像；GDAL、PROJ与rasterio的版本组合必须在构建环境验证，不直接承诺任意系统\file{pip install gdal}即可运行。首版可用SQLite存运行索引，本地文件系统存大型产物；数据卷、数据库与审核日志一起备份。版本冻结清单包含代码、数据、参数、公式、Dify工作流、模型服务名和提示词版本。

\subsection{启动与交付检查}
\begin{enumerate}
\item 按交付README装载固定版本镜像，配置数据卷与可选LLM连接；首次运行先检查证据文件哈希。
\item 启动核算API和界面，再启动已配置的Dify服务；探测健康接口并验证工具网络。
\item 导入冻结基线和参数登记，验证报告、地图、指标卡指向同一run。
\item 运行手算检查和一次真实基线复算，再按第\ref{sec:demo}节脚本完成演练。
\item 冻结\file{haidian-demo-v0.1}交付包，附实际测量结果、未完成事项及责任人。只有通过验收的范围才能写入发布说明。
\end{enumerate}

\subsection{Overleaf同步约定}
详细实施方案源文件为\file{implementation/main.tex}。

概述方案位于\file{gep_demo_plan/main.tex}。详细版使用XeLaTeX编译及Fandol中文字体，减少宿主系统字体依赖。编辑前拉取远端最新版本，编译和差异检查后只提交明确修改的文件，普通推送，不覆盖他人更新。本地代码仓库\file{docs/implementation.tex}保留同内容副本。

\section{参考依据与修订记录}
\subsection{可定位的参考资料}
以下标准与资料用作依据；公开产品说明已在2026年9月21日审阅时核对。数据获取保留回执；使用时冻结请求、资产版本和许可，不将目录可访问当成已验收。
\begin{enumerate}[label={[\arabic*]}]
\item GB/T 46869--2025《生态系统评估 陆域生态产品总值核算技术指南》。项目本地PDF，71页。核查位置：第4--5章和表1--3（PDF第9--13页）；附录A（PDF第21--46页）；表C.21（PDF第65页）；附录D、E（PDF第67--69页）。
\item NASA MOD17A3HGF.061年NPP产品（2024年度候选核算输入）及Earth Engine目录：\href{https://developers.google.com/earth-engine/datasets/catalog/MODIS_061_MOD17A3HGF}{产品说明、波段单位及质量层}。
\item Google/WRI Dynamic World V1：\href{https://developers.google.com/earth-engine/datasets/catalog/GOOGLE_DYNAMICWORLD_V1}{类别概率、标签与使用说明}。
\item ESA/Copernicus Sentinel-2协调地表反射率产品：\href{https://developers.google.com/earth-engine/datasets/catalog/COPERNICUS_S2_SR_HARMONIZED}{波段、质量标记与比例因子}。
\item Copernicus CDS：\href{https://cds.climate.copernicus.eu/datasets/reanalysis-era5-land-monthly-means?tab=overview}{ERA5-Land月平均产品}；ECMWF：\href{https://confluence.ecmwf.int/pages/viewpage.action?pageId=505384848}{累积量、月平均与变量说明}。
\item NASA SRTM：\href{https://developers.google.com/earth-engine/datasets/catalog/USGS_SRTMGL1_003}{30 m高程产品说明}。
\item geoBoundaries：\href{https://www.geoboundaries.org/}{公开行政边界目录}。具体海淀边界资产与适用年份需获取后核验。
\item Dify官方文档：\href{https://docs.dify.ai/en/self-host/deploy/quick-start/docker-compose}{Docker Compose部署与初始化}。部署时选定版本并保存对应配置。
\item Google/Google DeepMind AlphaEarth年度嵌入：\href{https://developers.google.com/earth-engine/datasets/catalog/GOOGLE_SATELLITE_EMBEDDING_V1_ANNUAL}{官方目录、年份与获取入口}。成果保留署名：The AlphaEarth Foundations Satellite Embedding dataset is produced by Google and Google DeepMind.
\item 海淀区统计局：\href{https://zyk.bjhd.gov.cn/sjkf/ndsj/202601/t20260104_4799659.shtml}{2025北京海淀统计年鉴}，官网注明最新数据截至2024年度。
\item 海淀区：\href{https://zyk.bjhd.gov.cn/sjkf/tjgb/202504/t20250423_4766489.shtml}{2024年国民经济和社会发展统计公报}；\href{https://zyk.bjhd.gov.cn/sjkf/tjgb/}{统计公报年度目录}。
\item 北京市水务局：\href{https://swj.beijing.gov.cn/zwgk/szygb/202504/P020250428362976968562.pdf}{北京市2024年水资源公报}。市级背景与总量校核，不替代海淀模型输入。
\item 北京市：\href{https://www.beijing.gov.cn/ywdt/zwzt/jjjxtfz/cg/sthjbh/gb/202505/P020250520406072608853.pdf}{2024年北京市生态环境状况公报}。
\item 北京市水务局：\href{https://swj.beijing.gov.cn/zwgk/stbcgb/202505/P020250526402026040857.pdf}{2024年北京市水土保持公报}。
\item NASA MOD16A2GF.061：\href{https://developers.google.com/earth-engine/datasets/catalog/MODIS_061_MOD16A2GF}{8日蒸散、比例因子与质量标记}。
\item Source Cooperative / Taylor Geospatial：\href{https://source.coop/tge-labs/aef}{AlphaEarth年度数据、索引及README}。包括空间筛选、VRT方向与解量化要求，查阅日期2026-09-21。
\item Ouyang Z.等（2020）：\href{https://doi.org/10.1073/pnas.1911439117}{Using gross ecosystem product (GEP) to value nature in decision making}，PNAS 117(25):14593--14601。
\item 国家统计局：\href{https://www.stats.gov.cn/hd/cjwtjd/202302/t20230207_1902270.html}{农业统计常见问题：播种面积与单产调查}；\href{https://www.stats.gov.cn/zs/tjws/tjdc/202301/t20230101_1903786.html}{抽样框、入样概率与加权估计}。
\item FAO：\href{https://www.fao.org/fishery/static/FAO_Training/FAO_Training/General/x6709e/x6709e16.htm}{Management for freshwater fish culture：养殖记录与产出}。用于区分投苗、收获、库存和净增长。
\item 欧盟JRC：\href{https://joint-research-centre.ec.europa.eu/photovoltaic-geographical-information-system-pvgis/using-pvgis-5/pvgis-5-user-manual_en}{PVGIS 5 user manual}。光伏额定功率、辐照及损失参数定义。
\item 生态环境部：\href{https://www.mee.gov.cn/xxgk2018/xxgk/xxgk01/202106/t20210618_839512.html}{排放源统计调查产排污核算方法和系数手册发布公告}，2021年第24号。使用时固定对应行业手册版本。
\item 海淀区生态环境局：\href{https://zyk.bjhd.gov.cn/ztzl/hjbh/hjbhxw/202501/t20250117_4754538_m.shtml}{2024年度主要污染物浓度发布}，2025-01-17。
\item Nowak D.J., Crane D.E., Stevens J.C.（2006）：\href{https://research.fs.usda.gov/treesearch/14743}{Air pollution removal by urban trees and shrubs in the United States}，Urban Forestry \& Urban Greening 4:115--123；DOI:10.1016/j.ufug.2006.01.007。
\item Natural Capital Project：\href{https://storage.googleapis.com/releases.naturalcapitalproject.org/invest-userguide/latest/en/annual_water_yield.html}{InVEST Annual Water Yield模型文档}，用于方法边界核对，执行时固定软件版本。
\item Natural Capital Project：\href{https://storage.googleapis.com/releases.naturalcapitalproject.org/invest-userguide/latest/en/sdr.html}{InVEST SDR模型文档}，区分避免侵蚀与避免泥沙输出。
\item 河南省发展改革委：\href{https://sanmenxia.zfcg.henan.gov.cn/cmsweb81e27e/henan/rootfiles/2024/10/10/30b0c7e39e1d4f029e0b7384127972e8.pdf}{2024年节能降碳基础能力建设项目包9招标文件}，豫财招标采购-2024-981，第49页采购需求。
\item 云南省统计局：\href{https://www.ccgp.gov.cn/cggg/dfgg/jzxcs/202604/t20260417_26423542.htm}{2026年生态产品总值核算技术服务项目竞争性磋商公告}，2026-04-17。
\item ModelScope官方Hub SDK：\href{https://github.com/modelscope/modelscope_hub}{modelscope-hub源码及使用说明}。本项目锁定0.4.3，使用HubApi、dataset仓库类型与逐文件回下载验证。
\end{enumerate}

\subsection{本次修订摘要}
本版围绕2024数据落盘、国标分项计算、人工补数和可追溯发布细化实施步骤。第1.2节解释三类价值与两项调整；第2.1节登记2024年度数据及其在智能体中的用途；第3章补充统计品类、作物识别适用条件和农林牧渔采集表。新增代码操作手册与产品交互流程，通过玉米计价和水产缺数两个案例说明输入、选择、审核和输出；GitHub保存代码环境，ModelScope在本地冻结核验后接收版本化数据包。
\end{document}
